\section{64 位入口先把加载器状态收敛为内核 ABI}

Stage 1 已经把 RSP 放到独立的 64 KiB 栈顶，用 \code{CALL} 进入固定地址
\code{OsKernelEntry}，并在 RDI 中传入 BootInfo。这里故意使用正常 System V
AMD64 调用形状：函数入口的 RSP 模 16 等于 8，编译器可以按普通 C++20 ABI
生成代码。入口函数自身只有一项职责：选择“无故障注入”并进入
\code{RunKernel}，不复制启动逻辑。

内核先重新初始化 COM1，然后逐字段验证 BootInfo、BSS 探针和 CR3。任何失败
只输出一次对应边界并停机。只有这三项都成立，才开始覆盖 Stage 1 留下的 GDT。
这种顺序很重要：如果 BootInfo 已损坏，内核不能拿其中的栈顶去构造 TSS；
如果 BSS 没有清零，全局描述符表和不可重入状态也没有语言层保证。

System V AMD64 ABI 区分调用者与被调用者保存寄存器。内核入口虽然由加载器
\code{CALL}，但协议规定它永不返回；若意外返回，Stage 1 会输出
\code{KERNEL\_RETURNED} 并停机。异常入口则不是普通调用，必须由单独汇编边界
转换后才能进入 C++。

\section{标识符在跨语言边界上也是 ABI}

C++ 标识符主要服务于类型检查和源码阅读，NASM 标签与链接脚本符号则直接进入
ELF 符号表。两边都区分大小写，\code{OsKernelEntry} 与
\code{osKernelEntry} 是两个完全不同的链接名字。早期项目把一部分变量写成
小驼峰，也把自研 C/汇编入口写成小写开头；代码量较小时尚可辨认，模块增多后，
函数、变量和链接数据在走读中很难一眼区分。

当前规则按语义角色拆开：

\begin{longtable}{L{0.25\textwidth}L{0.27\textwidth}L{0.38\textwidth}}
\toprule
\textbf{对象} & \textbf{形式} & \textbf{示例}\\
\midrule
\endfirsthead
\toprule
\textbf{对象} & \textbf{形式} & \textbf{示例}\\
\midrule
\endhead
类型与函数 & 大驼峰 & \code{PhysicalAddress}、\code{RunKernel}\\
局部变量与参数 & 小写蛇形 & \code{physical\_address}、\code{frame\_index}\\
私有数据成员 & 小写蛇形，可带尾部下划线 &
\code{free\_frame\_count\_}\\
项目常量 & 项目/模块/功能全大写 &
\code{OS\_KERNEL\_MEMORY\_PAGE\_SIZE\_BYTES}\\
自研汇编函数符号 & 大驼峰 & \code{OsKernelExceptionDispatch}\\
自研链接数据符号 & 小写蛇形 & \code{os\_kernel\_image\_start}\\
\bottomrule
\end{longtable}

成员尾部下划线解决参数遮蔽，\code{this->} 则说明一次访问确实读取对象状态；
二者用途不同。命名空间的每一层只放一个简短小写单词，例如
\code{os::kernel::memory}，需要更细归属时增加层级，不把多个模块压成一个
冗长名称。Clang-Tidy 的 \code{lower\_case} 风格仍允许下划线，不能单独证明
“每层一个单词”；因此同一 CTest 先用 Python 词法扫描要求每层匹配
\code{[a-z]+}，再用 Clang AST 区分变量、函数、成员和类型。词法扫描会屏蔽
注释、普通/原始字符串和字符字面量，匿名命名空间则按 C++ 语义保留。

\code{extern "C"} 只关闭 C++ 名字改编，不会替我们同步 NASM 或链接脚本。
因此一次符号重命名必须同时修改声明、定义、\code{global}/\code{extern}、
链接器 \code{ENTRY}、ELF 审计器和调试文档。项目自己拥有这些 ABI，所以函数
统一为 \code{OsKernel...}/\code{OsUser...}，数据统一为
\code{os\_kernel\_...}；只有 C 标准和编译器明确要求的
\code{main}、\code{memset}、\code{memcpy} 保留固定拼写。

文本正则无法判断一个单词究竟是类型、函数还是变量。本次迁移用 Clang AST 扫描
136 个 C++ 头文件和源文件；初始 7018 条编译单元诊断包含公共头文件的重复报告，
迁移后诊断为零。加入 buddy、类型缓存与 KVA 测试后，当前门禁扫描 150 个
文件，并在每次完整回归中重新执行同一检查，使命名规范从说明文字变成可失败
的工程契约。

\section{BSS 让大表和专用栈不膨胀磁盘文件}

链接脚本把代码、只读数据和可写/BSS 分成 \code{R E}、\code{R}、
\code{RW} 三个 \code{PT\_LOAD}。Stage 1 第二遍装载时复制
\code{p\_filesz}，再清零到 \code{p\_memsz}；内核入口只验证结果，不重复
清零。

IDT 占 4096 字节，三个 IST 栈合计 48 KiB，再加 GDT、TSS 和 panic 状态，
绝大多数初值都是零。把它们放入 BSS 后，生产内核文件约 58 KiB，而 RW
加载段的文件长度仍为零、内存长度约 52 KiB。若把零字节实存进 ELF，磁盘
读取和 CRC 都会为没有信息量的内容付出成本。

BSS 探针必须在任何描述符对象使用前保持零。它把 C++ 静态存储期的语言保证
连接到 ELF \code{p\_memsz} 语义：日志出现 \code{BSS\_ZEROED}，才说明
加载器真正为后续内核对象建立了确定初态。

\section{为什么 64 位模式没有删除 GDT}

8086 用段寄存器参与地址计算；80286 和 80386 把段值改成描述符选择子，引入
base、limit、type、DPL 和 present。长模式随后把大多数代码/数据段的 base
和 limit 忽略，把地址隔离主责交给页表，但保留 CS 的执行模式与特权语义，
也保留门描述符和 TSS。

这种历史兼容带来一个看似矛盾的事实：64 位内核采用平坦地址空间，却仍必须
拥有 GDT。原因不是还要用传统分段切内存，而是：

\begin{itemize}
  \item CS 必须引用 present 的 64 位代码段，L 位决定 64 位指令语义；
  \item IDT gate 必须给出目标代码段选择子；
  \item TSS 描述符仍位于 GDT，\code{LTR} 通过选择子找到它；
  \item v0.8 的 Ring 3 到 Ring 0 控制转移依赖特权级和 RSP0。
\end{itemize}

Stage 1 的 GDT 是模式切换工具，不是内核资产。v0.5 用内核 BSS 中的新表覆盖
GDTR，从依赖方向上切断内核对加载器私有内存的长期依赖。

\section{v0.5 五槽 GDT 的逐位结构}

\begin{longtable}{L{0.10\textwidth}L{0.15\textwidth}L{0.25\textwidth}L{0.36\textwidth}}
\toprule
\textbf{索引} & \textbf{选择子} & \textbf{描述符} & \textbf{关键属性}\\
\midrule
\endfirsthead
\toprule
\textbf{索引} & \textbf{选择子} & \textbf{描述符} & \textbf{关键属性}\\
\midrule
\endhead
0 & \code{0x00} & null & 全零，捕获误用空选择子\\
1 & \code{0x08} & Ring 0 code & P=1、DPL=0、execute/read、L=1、D=0\\
2 & \code{0x10} & Ring 0 data & P=1、DPL=0、read/write、L=0\\
3--4 & \code{0x18} & 64-bit TSS & type=9、P=1、完整 64 位 base\\
\bottomrule
\end{longtable}

选择子低两位是 RPL，第 2 位选择 GDT 或 LDT，其余高位是表索引。因此索引
1、2、3 左移三位分别得到 \code{0x08}、\code{0x10}、\code{0x18}。
TSS 描述符需要 16 字节：传统 8 字节只有位置保存 32 位 base，后一个槽保存
base[63:32]。构造器把 limit[19:0] 和 base[63:0] 分散写入硬件规定字段，
宿主单元测试再完整解码，避免“QEMU 恰好只用了低地址”掩盖高位错误。

GDTR 的 limit 是最后一个有效字节偏移，不是条目数或字节数。五槽表总长
40 字节，所以 limit 为 39。项目在加载后执行 \code{SGDT} 比较 base 和
limit，使“执行过 LGDT”升级为“处理器确实引用预期表”。

\section{LGDT 不会替我们刷新段缓存}

\code{LGDT} 只修改 GDTR，不自动重新解释当前 CS、SS 或数据段隐藏缓存。
内核先加载 GDTR，再把代码选择子和一个本地返回地址压栈，执行
\code{RETFQ}。远控制转移同时装载新的 CS 描述符和 RIP。到达新标签后才重载
DS、ES、FS、GS、SS，最后执行 \code{LTR}。

顺序不能反过来：TSS 选择子只有在新 GDT 可见后才有效；普通近跳转不能刷新
CS；只改变可见选择子而忽略隐藏描述符，会让日志看起来正确、处理器实际仍沿用
旧属性。v0.5 最终用 CS、SS 和 \code{STR} 回读分别得到
\code{0x08}、\code{0x10}、\code{0x18}。

这是 v0.5 的阶段形状。v0.8 在 Ring 0 数据段之后插入用户数据段与用户
代码段，TSS 随之移到索引 5--6、选择子 \code{0x28}；对应的七槽结构在
用户边界章节逐项展开。把两个阶段区分开，才能理解选择子变化是 GDT 演进，
而不是文档与二进制冲突。

\section{104 字节 TSS 不再承担硬件任务切换}

早期保护模式把寄存器现场也放进 TSS，可用硬件 task switch 切换任务。现代
操作系统通常自行保存上下文，因为软件调度更灵活、成本和副作用更可控。
长模式 TSS 因此缩减为特权栈和 I/O bitmap 等关键状态。

\begin{longtable}{L{0.14\textwidth}L{0.14\textwidth}L{0.26\textwidth}L{0.32\textwidth}}
\toprule
\textbf{偏移} & \textbf{宽度} & \textbf{字段} & \textbf{当前值}\\
\midrule
\endfirsthead
\toprule
\textbf{偏移} & \textbf{宽度} & \textbf{字段} & \textbf{当前值}\\
\midrule
\endhead
\code{0x04} & 64 & RSP0 & BootInfo 的内核初始栈顶\\
\code{0x0C} & 64 & RSP1 & 0\\
\code{0x14} & 64 & RSP2 & 0\\
\code{0x24} & 64 & IST1 & 双重故障栈顶\\
\code{0x2C} & 64 & IST2 & NMI 栈顶\\
\code{0x34} & 64 & IST3 & 机器检查栈顶\\
\code{0x3C..0x5C} & 64 & IST4..IST7 & 0\\
\code{0x66} & 16 & I/O bitmap offset & 104\\
\bottomrule
\end{longtable}

每个 IST 栈为 16 KiB，互不重叠，向低地址增长。I/O bitmap offset 等于结构
大小，已经超过 TSS 的 inclusive limit 103，表示当前 TSS 内没有权限位图。
执行 \code{LTR} 后，处理器会把 GDT 内描述符从 available type 9 改成 busy
type 11；验证程序若仍要求字节值为 9，反而误解了硬件副作用。

v0.6 把每个 IST 存储块扩为 20 KiB：低端 4 KiB 保持 guard，随后 16 KiB
才是可写栈。内核新页表明确跳过三个 guard 物理页；TSS 的 IST 指针仍指向可写
部分的高端。这样栈向下越界会遇到 not-present 页，而不是静默写入相邻对象。

\section{IDT 门描述符定义异常入口的硬件动作}

IDTR 保存 IDT 基址和界限。64 位门描述符包含 64 位处理函数地址、段选择子、
IST 索引、门类型、DPL 和 present 位。处理器根据向量索引门，检查权限，
必要时切换栈，再压入返回现场。

\begin{longtable}{L{0.19\textwidth}L{0.24\textwidth}L{0.43\textwidth}}
\toprule
\textbf{门字段} & \textbf{取值策略} & \textbf{作用}\\
\midrule
\endfirsthead
\toprule
\textbf{门字段} & \textbf{取值策略} & \textbf{作用}\\
\midrule
\endhead
offset & 异常桩地址 & 分三段编码完整 64 位入口\\
selector & Ring 0 代码段 & 进入内核执行环境\\
type & interrupt/trap gate & 决定入口是否自动清 IF\\
DPL & 通常 0 & 限制用户软件主动触发\\
IST & 关键异常专用栈 & 在当前栈损坏时仍可诊断\\
present & 1 & 允许处理器使用该门\\
\bottomrule
\end{longtable}

陷阱门保留 IF，适合调试类异常；中断门清 IF，避免普通硬件中断立刻嵌套。
选择是并发策略的一部分，不能对所有向量机械使用同一类型。

\section{TSS 在 64 位模式主要提供特权栈}

长模式不使用 TSS 做硬件任务切换，但保留 RSP0--RSP2 和七个 IST 指针。
从 Ring 3 进入 Ring 0 时处理器使用 RSP0；设置 IST 的门无条件切到对应栈。
双重故障、NMI 和机器检查适合使用独立 IST，避免普通内核栈溢出导致再次故障。

TSS 描述符占两个 GDT 槽，必须编码完整 64 位基址。加载 TR 前，TSS 内存、
GDT 描述符和专用栈都要映射并保持生命周期覆盖整个内核运行。

\section{异常向量具有不同错误码和恢复语义}

\begin{longtable}{L{0.12\textwidth}L{0.22\textwidth}L{0.17\textwidth}L{0.37\textwidth}}
\toprule
\textbf{向量} & \textbf{异常} & \textbf{错误码} & \textbf{初期策略}\\
\midrule
\endfirsthead
\toprule
\textbf{向量} & \textbf{异常} & \textbf{错误码} & \textbf{初期策略}\\
\midrule
\endhead
0 & 除法错误 & 无 & 记录 RIP 与寄存器后 panic\\
6 & 非法指令 & 无 & 反汇编故障地址并 panic\\
8 & 双重故障 & 固定错误码 & IST 栈上最小诊断并停机\\
13 & 一般保护 & 有 & 解码选择子相关位并 panic\\
14 & 页故障 & 有 & 结合 CR2 分类映射或非法访问\\
\bottomrule
\end{longtable}

某些异常属于 fault，保存 RIP 指向故障指令，修复原因后可以重试；trap 通常
保存下一指令；abort 不保证可恢复。处理程序必须知道类别，不能统一递增 RIP
跳过错误。

\section{统一异常帧让 C++ 分发器看到稳定布局}

处理器有时压错误码，有时不压；特权切换时还会多压 SS 与 RSP。每个汇编桩
补齐缺失字段并保存通用寄存器，把向量号与规范化错误码放到固定位置。C++ 只
接受指向统一 \code{InterruptFrame} 的指针。

结构字段偏移由汇编常量和 C++ \code{static_assert} 双向检查。新增寄存器或
改变压栈顺序必须同时更新弹栈路径；\code{iretq} 前残留一个八字节字段都会把
后续值错当 RIP 或 CS。

\section{v0.5 的 32 个桩怎样形成 160 字节帧}

CPU 对向量 8、10、11、12、13、14、17、21、29、30 压入错误码，其余异常
不压。NASM 用两种项目内宏生成桩；宏只减少机械重复，不隐藏硬件策略：

\begin{lstlisting}[language={[x86masm]Assembler}]
; 无硬件错误码
push qword 0
push qword vector
jmp OsKernelExceptionDispatch

; 有硬件错误码：CPU 已经压入 error_code
push qword vector
jmp OsKernelExceptionDispatch
\end{lstlisting}

两条路径到达公共入口时，栈顶永远是 vector，其后永远是 error code，再后是
CPU 保存的 RIP、CS、RFLAGS。公共入口执行 \code{CLD}，因为 C++ ABI 假定
方向标志清零；被打断代码若设置过 DF，\code{IRETQ} 最终会从保存的 RFLAGS
恢复原值。

寄存器依次压入 RAX、RBX、RCX、RDX、RBP、RSI、RDI、R8 到 R15。栈向低地址
增长，所以 C++ 看到的结构顺序正好反过来：

\begin{longtable}{L{0.16\textwidth}L{0.23\textwidth}L{0.47\textwidth}}
\toprule
\textbf{帧内偏移} & \textbf{字段} & \textbf{来源}\\
\midrule
\endfirsthead
\toprule
\textbf{帧内偏移} & \textbf{字段} & \textbf{来源}\\
\midrule
\endhead
\code{0x00..0x38} & R15..R8 & 公共桩保存\\
\code{0x40..0x70} & RDI..RAX & 公共桩保存\\
\code{0x78} & vector & 每向量桩\\
\code{0x80} & error code & CPU 或零占位\\
\code{0x88} & RIP & CPU\\
\code{0x90} & CS & CPU\\
\code{0x98} & RFLAGS & CPU\\
\bottomrule
\end{longtable}

共 20 个 64 位字段，即 160 字节。v0.5 的异常来自 Ring 0，没有特权级切换，
所以 CPU 不另压旧 RSP 和 SS。v0.8 的 Ring 3 入口在公共结构之后真实增加
旧 RSP 与 SS，形成 176 字节 \code{UserPrivilegeFrame}；公共前缀保持不变。

C++ 调用前不能假设异常发生时 RSP 恰好满足 ABI。公共桩先把帧地址放入 RDI，
再把 RSP 向下按 16 字节对齐，预留一槽保存原帧地址后调用分发器。若分发器
返回，就从预留槽恢复原 RSP，逆序弹出寄存器，丢弃 vector/error 两槽并执行
\code{IRETQ}。普通 \code{RET} 无法恢复 CS 和 RFLAGS。

\section{全部异常先被分类，再决定是否恢复}

\begin{longtable}{L{0.08\textwidth}L{0.14\textwidth}L{0.17\textwidth}L{0.20\textwidth}L{0.27\textwidth}}
\toprule
\textbf{向量} & \textbf{名称} & \textbf{类别} & \textbf{错误码} & \textbf{v0.5 策略}\\
\midrule
\endfirsthead
\toprule
\textbf{向量} & \textbf{名称} & \textbf{类别} & \textbf{错误码} & \textbf{v0.5 策略}\\
\midrule
\endhead
0 & \#DE & fault & 无 & panic\\
1 & \#DB & fault/trap & 无 & panic\\
2 & NMI & interrupt & 无 & IST2 后 panic\\
3 & \#BP & trap & 无 & 仅启动自检返回\\
4 & \#OF & trap & 无 & panic\\
5 & \#BR & fault & 无 & panic\\
6 & \#UD & fault & 无 & 注入测试后 panic\\
7 & \#NM & fault & 无 & panic\\
8 & \#DF & abort & 固定 0 & IST1 后 panic\\
10--14 & \#TS..\#PF & fault & 有 & panic\\
16 & \#MF & fault & 无 & panic\\
17 & \#AC & fault & 有 & panic\\
18 & \#MC & abort & 无 & IST3 后 panic\\
19--20 & \#XM/\#VE & fault & 无 & panic\\
21 & \#CP & fault & 有 & panic\\
28 & \#HV & fault & 无 & panic\\
29--30 & \#VC/\#SX & fault & 有 & panic\\
\bottomrule
\end{longtable}

表中未列出的保留向量也有独立桩并进入 panic。只允许 breakpoint 返回不是
因为其他 fault 永远不能恢复，而是 v0.5 还没有足够策略证明恢复安全。例如
页不存在可以在 v0.6 分配并映射后重试；非法指令通常意味着代码或 CPU 能力
契约错误，跳过一个字节既不知道指令长度，也会掩盖破坏。

\section{INT3、UD2 和页故障组成三角验收}

单一正常启动不能证明异常系统。v0.5 用三条具有不同语义的真实指令路径：

\begin{enumerate}
  \item 正常内核执行 \code{INT3}。CPU 保存下一条 RIP，分发器确认
        vector=3、error=0，输出 \code{BREAKPOINT\_HANDLED} 后返回；
        只有随后出现 \code{EXCEPTION\_SELF\_TEST\_READY} 和
        \code{READY}，才证明 \code{IRETQ} 往返正确。
  \item 非法指令内核执行 \code{UD2}。它是架构保证始终触发 \#UD 的两字节
        指令，避免用“碰巧非法”的字节序列。panic 必须报告 vector=6、
        error=0，禁止输出 \code{READY}。
  \item 页故障内核读取 \code{0x04000000}。Stage 1 只映射低 64 MiB，
        因此该值是第一个确定未映射地址。panic 必须报告 vector=14、
        error=0、CR2 同值，并禁止 \code{READY}。
\end{enumerate}

页故障错误码零逐位表示 P=0、W/R=0、U/S=0、RSVD=0、I/D=0：supervisor
读取 not-present 的数据页。这比只看到“向量 14”提供更多证据；若访问意外
变成写、用户态或权限违反，错误码会改变，测试能及时暴露。

两份故障内核不复制描述符和异常实现，只用最薄入口选择注入模式。否则测试代码
可能通过，而实际生产入口链接的是另一套未经验证的处理器。

\section{panic 的可信计算基}

panic 可能由坏页表、坏栈、坏堆或坏锁引起。让 panic 调用通用格式化器、分配
字符串或获取日志锁，会扩大它必须信任的系统。v0.5 把可信集合压到：

\begin{itemize}
  \item 已在入口独立初始化的 COM1 端口；
  \item 当前统一异常帧；
  \item CR2 等直接可读架构寄存器；
  \item 一个零初始化的 64 位不可重入状态。
\end{itemize}

panic 首先 \code{CLI}，发现状态已 active 就直接 HLT，不再次打印。首次进入
依次输出向量、错误码、RIP、CS、RFLAGS；只有页故障读取 CR2。每个字段使用
固定 16 位十六进制宽度，不需要除法、区域设置或动态缓冲。最后输出一次
\code{PANIC} 并进入 \code{CLI; HLT} 循环。

RIP 会随链接变化，RFLAGS 也可能包含处理器为 fault 设置的 RF 等位，因此
QEMU 协议要求字段存在，却不硬编码这两个不稳定数值。向量、错误码、页故障
地址和终止标记才是阶段稳定主键。这正是“日志丰富但不刷屏”的工程边界。

\section{从位字段单元测试到整机失败路径}

描述符纯编码使用宿主 C++ 测试，不必每次启动 QEMU。确定样例把 TSS base、
limit、type 和 IDT offset、selector、IST、attributes 逐字段解码。固定种子
\code{0xD35C71A05EED6405} 再生成 4096 个完整 64 位地址，分别经过 TSS
descriptor 和 IDT gate 编码解码，共执行 8192 条往返性质断言。

集成审计读取真实 ELF，要求 \code{OsKernelLoadGdtAndTss}、
\code{OsKernelExceptionDispatch}、\code{OsKernelDispatchException}、
桩表和向量 0..31 的全部符号存在，同时拒绝任何未解析运行时符号。

QEMU 则验证纯逻辑无法证明的事实：处理器接受 GDT/TSS/IDT，TR 真正装载，
INT3 能返回，UD2 和页故障进入正确向量，panic 后不会继续成功路径。v0.5
完成时 36 项 CTest 全部通过，其中 17 项带失败路径标签；既有复位、ATA、
长模式和 ELF 回归也必须同时通过。

v0.6 在这条回归上增加七项测试，完整集合为 43 项。新增 QEMU 写保护镜像不
复制异常实现，只让同一个生产内核入口选择 \code{WriteProtection} 注入；
因此错误码 3、CR2 和 PANIC 同时验证页表、CR0.WP、IDT、统一帧与串口路径。

\section{页故障错误码提供访问分类}

页故障错误码的 P 位区分不存在与权限违反，W/R 区分写与读，U/S 区分用户与
内核，RSVD 表示页表保留位错误，I/D 表示取指。CR2 保存故障线性地址。

用户访问不存在页可能触发按需分配，用户写只读共享页可能触发写时复制，用户
访问内核页应终止进程；内核访问非法地址通常是内核缺陷。处理器给出的位只描述
机制原因，内核还要结合虚拟区间元数据决定策略。

\section{内存图在加载器排序，在内核验证}

当前平台后端是 QEMU PC \code{fw\_cfg}。Stage 1 从
\code{etc/e820} 读取 20 字节原始项，扩展成：

\begin{lstlisting}[language=C++]
struct PhysicalMemoryMapEntry final {
    uint64_t base_address;
    uint64_t length_bytes;
    uint32_t type;
    uint32_t attributes;
};
\end{lstlisting}

结构精确为 24 字节。base 和 length 必须保持 64 位，因为物理地址可能远高于
当前管理窗口；type/attributes 是外部格式的 32 位字段，不能为了“用大的”
擅自改变 ABI。Stage 1 在交接前按 base 排序，BootInfo v2 提供
\code{0x18000}、条目数和条目宽度。

内核不重新排序，也不尝试猜测重叠优先级。它把“已排序、互不重叠”视为交接
契约并独立验证：

\begin{enumerate}[label=\arabic*.]
  \item 指针非空，数量在 1 到 128，管理上限非零；
  \item 每项 length 非零，base+length 不发生 64 位溢出；
  \item 当前 base 不小于前一项 base，也不小于前一项 end；
  \item 所有描述长度求和、type 1 可用长度求和均不溢出；
  \item 由最高 type 1 RAM 页、处理器物理地址宽度与 direct-map 容量共同
        确定的管理范围内至少存在完整可用页。
\end{enumerate}

乱序和重叠使用不同状态返回，因为前者通常是平台规范化缺失，后者表示两个
所有者对同一物理地址作了冲突声明。内核若自行切分并选择“保留优先”，反而会
把上游错误悄悄变成某种策略；早期系统选择明确拒绝，现场更容易解释。

汇总输出先写入局部候选对象，全部成功后才覆盖调用者。这样某个高地址条目在
最后溢出时，调用者不会误用前面已经累加的半份统计。日志分别打印
\code{MEMORY\_DESCRIBED\_BYTES}、\code{MEMORY\_USABLE\_BYTES} 和
\code{MEMORY\_MANAGED\_BYTES}；高端保留 MMIO 窗口只进入第一个值，不能
冒充 RAM。

\section{2-bit 页帧状态把所有权错误变成显式状态}

v0.6 最初固定管理低 64 MiB；v1.1 已把这个教学脚手架替换为运行期容量。内核
先读取 \code{CPUID.80000008H:EAX[7:0]} 的物理地址宽度，再取三个上限
的最小值：

\begin{enumerate}[label=\arabic*.]
  \item E820 中最高 type 1 RAM 完整页的末端；
  \item 处理器能够表达的物理地址上限；
  \item 四级页表为高半区 direct-map 保留的 64 TiB。
\end{enumerate}

这里必须使用“最高可用 RAM”而不是“最高被描述地址”。64 MiB QEMU 配置会
同时描述一个延伸到 1 TiB 的保留窗口；保留洞要阻止分配，却不能让 64 MiB
机器为 1 TiB PFN 建表。

每帧不用一位，而用两位表示四种状态：

\begin{longtable}{L{0.18\textwidth}L{0.25\textwidth}L{0.43\textwidth}}
\toprule
\textbf{状态} & \textbf{来源} & \textbf{允许的迁移}\\
\midrule
\endfirsthead
\toprule
\textbf{状态} & \textbf{来源} & \textbf{允许的迁移}\\
\midrule
\endhead
unavailable & 默认、非 type 1、碎片页 & 不参与普通分配\\
free & 完整落入可用区域的 4 KiB 帧 & allocate 或 reserve\\
allocated & 动态分配成功 & 只能 release 回 free\\
reserved & 平台、Kernel、初始栈 & 普通接口不释放\\
\bottomrule
\end{longtable}

设管理上界为 \(L\)，页大小为 \(P=4096\)，状态存储为：
\[
  N_{\mathrm{frame}}=\frac{L}{P},\qquad
  S_{\mathrm{state}}=
  \left\lceil\frac{N_{\mathrm{frame}}}{4}\right\rceil.
\]
64 MiB 连续 RAM 仍只需 4096 字节。64 GiB QEMU 机器在 3--4 GiB 有一段
物理地址洞，最高 RAM 末端为 65 GiB，因此状态存储需要
\code{0x410000} 字节；洞中的 PFN 保持 unavailable，实际可用 RAM 仍精确为
64 GiB。
这里使用 \code{uint8\_t} 数组并不是违背内部大位宽规范；密集状态的每字节
恰好保存四帧，扩大为 64 位会把元数据增加八倍。对外地址、计数、容量和索引
仍全部使用 64 位。

初始化先把所需状态字节清零，对应 unavailable，再把可用区间按 4 KiB 向内
收缩。若区域起点不是页边界，向上对齐；终点向下对齐。只有整个页都被平台声明
为可用时才进入 free。连续四帧使用预计算的 \code{0x55} 状态字节批量提交，
避免 64 GiB 启动执行一千六百万次逐帧 read-modify-write；边界不足四帧时
仍逐项更新，不会把相邻洞误标为 free。全部扫描后仍没有 free 帧，分配器把
managed count 恢复为零再返回 \code{NoUsableFrames}，不会留下“看似初始化”
的对象。

\section{启动保留采用先检查、后提交}

状态数组本身不能继续放在固定 BSS。内核先根据上式计算并向页对齐所需大小，
再在 Stage 1 可访问的低 64 MiB RAM 中搜索连续区间；搜索器逐个跳过重叠
保留区，并在每次跳转后重新满足 4 KiB 对齐。内核随后依次保留四类范围：

\begin{itemize}
  \item \code{[0,1 MiB)}：复位、固件、Stage 1、旧页表、BootInfo、内存图、
        VGA 传统窗口等低端平台资产；
  \item 链接器符号 \code{os_kernel_image_start..os_kernel_image_end}：精确覆盖
        text、rodata、data、BSS、IDT 与 IST；
  \item BootInfo 给出的 64 KiB 初始栈。
  \item 动态选出的页帧状态数组与 buddy 双位图元数据。
\end{itemize}

任意起点和终点先向页边界扩张，保证部分接触的帧也被保留。
\code{ReserveRange} 第一遍只读状态；只要遇到 allocated 就立即失败。
第二遍才把 free 改为 reserved 并更新计数。若边扫描边修改，范围前半段可能
已经 reserved，后半段才发现冲突，调用者得到失败却留下不可恢复的部分副作用。

早期线性模型从 next-search 索引开始环形搜索，成功后把状态改为 allocated
并把下一搜索点移到后一帧。正式启动在产生首个动态页之前启用 buddy，此后
同一接口改走 order 0 位图，next-search 不再是权威空闲集合。两种阶段都验证
页对齐、管理范围和 allocated 状态；重复释放、释放 reserved 或 unavailable
都返回明确错误。

\code{AllocateInRange(begin,end)} 在同一状态机上增加页对齐的半开区间约束。
启动新 CR3 前用它把页表帧限制在低 64 MiB 身份映射中；高内存验收则把
begin 设为 4 GiB+4 KiB。范围接口没有复制第二套所有权规则，因此普通分配、
受限分配和释放共享重复释放与统计不变量。

单元测试专门构造“第一帧 free、第二帧 allocated”的两页保留请求，确认失败
后 free、allocated、reserved 三个计数完全不变。这条测试验证事务边界，而
不只是某个返回枚举。

\section{伙伴系统把连续页变成可证明的块}

\subsection{为什么逐页所有权还不足以支持连续申请}

每帧 2 bit 状态已经能回答“这个 4 KiB 页是否可以交给调用者”，却不能直接
回答“从这里开始的八页是不是一个可交付、可原样收回的连续对象”。在状态数组
上向后数八个 free 可以实现一次成功申请，但释放后还要重新发现相邻区间；
更危险的是，调用者若把八页申请中的第一页按单页释放，逐页状态本身无法知道
原所有权边界。

连续页并不是为了让普通小对象跳过 heap。它服务于页表批量后备、某些设备缓冲、
将来可能采用的物理大页，以及需要一次保留若干页的内核资源。KVA 与动态内核栈
通常只要求\emph{虚拟}连续，物理页可以离散；但物理层仍必须有明确提供连续块
的能力，且不能让这种能力破坏现有单页路径。

早期系统常在两种极端之间取舍：顺序表能精确表达任意长度，却需要搜索和合并
相邻区间；固定大小池查找很快，却只能服务一个尺寸。伙伴方法把尺寸限制为
二的幂，换取由地址位直接确定的唯一合并对象。这种限制带来内部碎片，但让
分裂、释放和局部一致性都能用短而确定的规则描述，适合先建立可信页资源地基。

\subsection{order 是页号二进制结构，不是随意标签}

order \(k\) 的块含 \(2^k\) 页，字节数为：

\[
  B(k)=2^k\times4096=2^{k+12}.
\]

块首 PFN 必须按 \(2^k\) 对齐。若首 PFN 为 \(p\)，同阶伙伴可写成：

\[
  \operatorname{buddy}(p,k)=p\oplus2^k.
\]

实现把 PFN 先右移 \(k\) 位得到块索引，所以等价代码是
\code{buddy\_block\_index = block\_index XOR 1}。这不是一种搜索启发式；
第 \(k\) 位恰好表示当前块处于同一父块的左半还是右半。只有这两个半块都空闲
时，清掉该位才得到唯一父块。

\begin{center}
\begin{tabular}{llll}
\toprule
order & 页数 & 字节数 & 典型含义\\
\midrule
0 & 1 & 4 KiB & 普通页表页、用户页\\
1 & 2 & 8 KiB & 小型连续后备\\
3 & 8 & 32 KiB & 本项目 QEMU 连续块自检\\
9 & 512 & 2 MiB & 与 x86-64 2 MiB 大页同尺寸\\
24 & 16777216 & 64 GiB & 64 GiB 页规模的理论最高阶\\
\bottomrule
\end{tabular}
\end{center}

“同尺寸”不代表当前 order 9 申请会自动安装成 2 MiB PDE。buddy 只交付物理
连续性；页表层是否使用 PS 位、虚拟地址是否对齐、权限是否一致仍由映射契约
决定。

\subsection{为什么保留 2-bit 状态并再加两族位图}

只用一张 free buddy 位图可以完成分裂和合并，却不能证明释放参数正是原申请。
例如 order 3 块首也满足 order 2 对齐；若释放接口盲信调用者传回的 order，
它可能释放前四页而让后四页永远悬空。于是每阶保存两张位图：

\begin{itemize}
  \item free 位图表示这一阶中哪些完整块是空闲集合成员；
  \item allocated 位图表示哪些块首曾以这一阶交付且尚未释放；
  \item 原 2-bit/PFN 状态继续区分 unavailable、reserved、free 与 allocated。
\end{itemize}

free 位图回答算法集合，allocated 位图回答块身份，2-bit 状态回答逐页硬件与
启动所有权。三者职责相邻但不可互相冒充。正式内核没有另一个可以绕过 buddy
的页分配器：旧 \code{Allocate}/\code{AllocateInRange}/\code{Release}
在启用后统一解释为 order 0，页表和用户地址空间无需改写公开单页契约。

设最高 PFN 范围含 \(N\) 帧，一族位图大小为：

\[
\begin{aligned}
  F(N) &=
  \sum_{k=0}^{\lfloor\log_2N\rfloor}
  \left\lceil
    \frac{\left\lceil N/2^k\right\rceil}{8}
  \right\rceil,\\
  M_{\mathrm{buddy}}(N)&=2F(N).
\end{aligned}
\]

块数随阶数形成 \(N+N/2+N/4+\cdots\)，所以两族总位数小于约 \(4N\)。
若为每个 PFN 保存一个 64 位 next 指针，64 GiB 就要 128 MiB；双位图的
同规格精确值只有 8388612 字节。

\begin{longtable}{L{0.18\textwidth}L{0.16\textwidth}L{0.19\textwidth}L{0.19\textwidth}}
\toprule
\textbf{PFN 覆盖} & \textbf{页数} & \textbf{buddy 精确字节} &
\textbf{页对齐字节}\\
\midrule
\endfirsthead
\toprule
\textbf{PFN 覆盖} & \textbf{页数} & \textbf{buddy 精确字节} &
\textbf{页对齐字节}\\
\midrule
\endhead
64 MiB & 16384 & 8200 & \code{0x3000}\\
64 GiB & 16777216 & 8388612 & \code{0x801000}\\
65 GiB 最高 PFN & 17039360 & 8519694 & \code{0x821000}\\
\bottomrule
\end{longtable}

QEMU 主规格报告 64 GiB 可用 RAM，但 3--4 GiB 洞把最高可用地址推到
65 GiB，所以线性 PFN 元数据按第三行计算。对应 2-bit 状态为
\code{0x410000} 字节，两类元数据合计 \code{0xC31000}，约 12.2 MiB。
这是按最高 PFN 换取简单索引的明确代价；极端稀疏物理地址空间以后需要
分段 vmemmap，而不是假称当前布局没有开销。

\subsection{启动时先冻结所有保留，再建立最大块分解}

页状态区和 buddy 区分别按 4 KiB 对齐，再合并成一个连续启动元数据事务。
搜索仍发生在 Stage 1 身份映射可访问的低 64 MiB，跳过低端平台、Kernel 和
初始栈。初始化顺序必须保持：

\begin{enumerate}
  \item 从 E820 建立 2-bit 页状态；
  \item 保留低端平台、Kernel、初始栈和完整元数据区；
  \item 确认 allocated 页计数仍为零；
  \item 扫描每一段连续 free PFN，建立 buddy；
  \item 此后才允许页表、heap 和用户页开始申请。
\end{enumerate}

第三步阻止一种不可恢复的猜测：若旧线性分配器已经交付若干单页，buddy 不知道
它们是否本来属于更大的调用者请求。当前启动流程不猜，而是返回
\code{ExistingAllocationsPreventBuddyInitialization}。

每段 free PFN 使用“满足当前首地址对齐的最大阶”贪心分解。设区间为
\([s,e)\)，先取不超过 \(e-s\) 的最高阶，再把阶降低到 \(s\) 也满足对齐；
登记该块后令 \(s\) 前进块大小并重复。这样 E820 洞、reserved 页和不完整尾部
自然切断块，同时不会留下本可直接合并的两个同阶伙伴。

buddy 建立后 \code{ReserveRange} 返回
\code{ReservationAfterBuddyInitialization}。运行期内存下线不是“把几页改成
reserved”这么简单；它要从空闲块拆出范围、迁移活动页并支持失败回滚，因此
必须作为以后独立协议实现。

\subsection{范围申请先选目标，再提交分裂}

\code{AllocateBlockInRange(order,begin,end)} 要求 begin/end 页对齐，并保证
最终目标块完整落入半开区间。算法从请求阶向上寻找最低存在候选的源阶。某个
较大源块可以只与请求范围部分相交，只要其中存在一个按请求阶对齐且完整位于
范围内的目标块。

找到源块并不立即改位图。实现先在局部变量中重放每一级“选择左半还是右半”，
确认最终正好到达目标 PFN，随后一次提交：

\begin{lstlisting}[language=C++]
remove free[source_order][source_index]
while current_order > requested_order:
    current_order -= 1
    keep the half containing target
    insert the other half into free[current_order]
insert allocated[requested_order][target_index]
mark every target page allocated
commit counters and output
\end{lstlisting}

假设从 order 3 的八页块申请 order 0 的一页，位图中会依次留下 order 2 的
四页伙伴、order 1 的两页伙伴和 order 0 的一页伙伴。源块不再同时存在；
任一 PFN 在任何时刻只能被一个已登记 free/allocated 块覆盖。

若总 free 页不少但碎片无法形成请求阶，结果仍是 \code{OutOfMemory}。这是
连续性契约的真实失败，不应退回八个离散页并假装成功。无效 order、范围未对齐、
越界和溢出同样在修改输出之前返回，调用者可以安全地进行作用域回滚。

\subsection{释放必须证明地址、阶和页状态三者一致}

\code{ReleaseBlock} 先验证页对齐、阶对齐、管理边界，再查精确
\code{allocated[order][index]}：

\begin{itemize}
  \item 同一物理位置被另一阶活动块覆盖，返回
        \code{AllocationOrderMismatch}；
  \item 没有活动块覆盖，返回 \code{FrameNotAllocated}，包括重复释放和
        reserved 页；
  \item 活动位存在但块内某页不是 allocated，返回 \code{CorruptedState}。
\end{itemize}

验证成功后才清活动位、把整块页状态恢复为 free。当前块此时尚未插入 free
位图；循环只查询并移除 XOR 伙伴，每成功一次就把块索引右移一位、order 加一。
最后只插入唯一最高合并结果。这一顺序避免“当前块和父块短暂同时可见”，也让
空闲块计数能够逐阶守恒。

\subsection{完整校验为何不需要再申请内存}

分配器不能为了验证自己再调用自己。因而 \code{ValidateBuddy} 逐阶扫描现有
位图，并利用块内页状态作交叉证据：

\begin{enumerate}
  \item 位图尾部超出完整块数的 bit 必须为零；
  \item 同块不能同时属于 free 与 allocated；
  \item 任意已登记块不能再有 free 或 allocated 祖先；
  \item 同阶 XOR 伙伴若都 free，说明遗漏了一次合并；
  \item free/allocated 块内每页必须具有对应 2-bit 状态；
  \item 各块按 \(2^k\) 加权后的页数必须等于全局页统计；
  \item 累计成功数减释放数必须等于活动块数。
\end{enumerate}

64 GiB 下这项检查会遍历覆盖最高 PFN 的元数据，所以 QEMU 宿主时间明显长于
64 MiB。日志中的 \code{[QEMU][T+......ms]} 让这段成本可见；不能通过只管理
64 MiB 或删除检查来制造更快的假结果。

\subsection{三层模型和真实执行各自证明什么}

单元测试以 64 页验证 36 字节精确位图、缺失/过小存储、已有活动页、最大块
分解、错阶、错位、重复释放和保留冻结。集成测试构造 1024 页并在中间放置
256 页 E820 洞，要求 order 5 块完整落入指定高地址区间，随后混合 order
0/3/5/6 并乱序回收。

固定种子 \code{0x425544445936344D} 执行 100000 步随机申请、释放、耗尽和
重复释放。独立布尔数组逐页记录参考所有权，每 257 步运行一次完整校验，失败
会报告种子和迭代号。

QEMU 不是再做一遍纯算法测试。64 MiB 与 64 GiB 都申请 order 3 的八页块，
通过 direct-map 查询首尾映射、写入
\code{0x4255444459464952} 与 \code{0x42554444594C4153}，读回后释放。
64 GiB 还要求物理地址高于 4 GiB、最大 order 至少 24。只有进入前后的页统计
与活动块数恢复基线，并通过完整校验，才输出
\code{BUDDY\_SELF\_TEST\_PASSED}。

日志中的 \code{BUDDY\_ACTIVE\_BLOCKS} 不应为零：正式页表、heap 后备页和
写保护测试页仍由内核持有。正确问题是“一段生命周期之后是否回到进入前”，
而不是“运行中的内核是否拥有零页”。当前单 BSP 路径由调用顺序串行化；
v1.2 已建立 Thread 与三类锁的调用边界，但尚未宣称分配器支持 SMP 并发。
未来开放并发调用时必须由外层物理内存锁保护；IRQ、NMI 和 panic 仍不得进入
通用分配路径。

\section{从 Stage 1 大页平滑切换到内核 4 KiB 页表}

构造新页表时处理器仍运行在 Stage 1 的 CR3 下。旧表把低 64 MiB 通过
32 个 2 MiB 大页恒等映射，所以新分配的页表帧可以直接把物理地址解释为指针
并清零。但这个等式只是启动前置条件，不是运行期 C++ 地址模型。
\code{PageTableMemoryAccess} 因而显式保存“物理内存通过哪个虚拟基址访问”
和“页表帧本轮最多从哪里分配”：

\begin{longtable}{L{0.28\textwidth}L{0.25\textwidth}L{0.35\textwidth}}
\toprule
\textbf{阶段} & \textbf{访问基址} & \textbf{页表帧分配上限}\\
\midrule
\endfirsthead
\toprule
\textbf{阶段} & \textbf{访问基址} & \textbf{页表帧分配上限}\\
\midrule
\endhead
旧 CR3 下构造 & 0，依赖身份映射 & 低 64 MiB\\
新 CR3 激活后 & \code{FFFF888000000000} & 全部受管 RAM\\
\bottomrule
\end{longtable}

内核先分配根 PML4，再对低 64 MiB 的启动访问逐页建立带精细权限的身份映射。
随后遍历 E820 type 1 RAM，在
\code{0xFFFF888000000000+physical} 建立 direct-map。未声明为 RAM 的
保留区与 MMIO 洞保持 not-present；LAPIC 继续使用独立 PCD 映射。

映射完成、NX 和 WP 打开之后才执行 \code{mov cr3,new_root}。如果先切 CR3
再补当前 RIP、RSP 或 IDT 所需映射，下一条取指或压栈就会故障；异常入口本身
也可能未映射，最终表现为 triple fault。先构造完整候选、后一次激活，是页表
版的两遍 ELF 副作用边界。CR3 读回成功后，管理器原子切换访问基址，页表
遍历、用户页清零和 ELF 段复制从此都不再把物理地址直接转换成指针。

\section{direct-map 为什么同时使用 2 MiB 与 4 KiB 页}

若把 64 GiB RAM 全部用 4 KiB PTE 映射，叶项数量为：
\[
 \frac{64\ \mathrm{GiB}}{4\ \mathrm{KiB}}=16777216,
\]
仅叶 PT 就需要约 32 MiB。使用 2 MiB PDE 大页时只需：
\[
 \frac{64\ \mathrm{GiB}}{2\ \mathrm{MiB}}=32768
\]
个叶项。大页显著减少页表内存，并让一个 TLB 项覆盖更大范围。

内核不能简单把每个 E820 条目向外取整到 2 MiB：那会把条目两侧的保留页也
映射成普通 RAM。算法先把起点向内对齐到 4 KiB，再用 4 KiB 页走到下一个
2 MiB 边界；中间完整区间使用大页，尾部回到 4 KiB。因而页粒度是每段的
局部选择，所有权边界仍由 E820 决定。

2 MiB PDE 设置 bit 7 的 PS，并把物理基址放在 bit 21..51；查询任意内部
地址时必须把虚拟地址低 21 位加回基址。普通用户页、Kernel 段和 guard page
仍使用 4 KiB PTE，因为它们需要细粒度权限或单页 not-present 语义。

\section{页表遍历必须验证每一级}

48 位 canonical 地址把 bit 47 符号扩展到 bit 63。若高 16 位既不是
\code{0x0000} 也不是 \code{0xFFFF} 对应的扩展，映射器在写表前就返回
\code{InvalidVirtualAddress}。四级索引为：
\[
\begin{aligned}
i_4&=(v\gg39)\mathbin{\&}0x1FF, &
i_3&=(v\gg30)\mathbin{\&}0x1FF,\\
i_2&=(v\gg21)\mathbin{\&}0x1FF, &
i_1&=(v\gg12)\mathbin{\&}0x1FF.
\end{aligned}
\]

中间表不存在时分配并清零。PML4E 与 PDPTE 出现 PS 仍属于当前实现不支持的
1 GiB 页并明确拒绝；PDE 的 PS=1 则是合法 2 MiB 叶，查询时使用
\code{0x000FFFFFFFE00000} 取物理基址。普通路径用
\code{0x000FFFFFFFFFF000} 取得下一级表。叶项写入前拒绝重复映射。
4 KiB 或 2 MiB 映射分别验证对应对齐，物理地址还不能越过处理器、页表字段
和 direct-map 访问环境的共同范围。

取消映射把叶项清零并执行 \code{INVLPG}，当前不回收空的 PT/PD/PDPT。
这是明确的早期边界：启动只映射固定集合，不会反复创建销毁地址空间。进程
阶段必须补充中间表引用计数或空表扫描，才能把页表帧安全归还分配器。

\section{权限在页表层级上累积}

任一上级项 RW=0 都会限制下级写，任一上级 US=0 都会限制用户访问。当前中间
表保持 RW，叶项决定具体内核页是否可写；用户映射请求出现时，映射器会把沿途
parent 的 US 置一，但 supervisor 叶项仍因自身 US=0 对用户不可见。

链接脚本导出 image、text、rodata、writable 的起止符号。内核按页应用：

\begin{longtable}{L{0.23\textwidth}L{0.18\textwidth}L{0.45\textwidth}}
\toprule
\textbf{区域} & \textbf{叶权限} & \textbf{理由}\\
\midrule
\endfirsthead
\toprule
\textbf{区域} & \textbf{叶权限} & \textbf{理由}\\
\midrule
\endhead
text & R/X & 指令可取，Ring 0 不可写\\
rodata & R/NX & 常量不可写，也不能作为代码执行\\
data/BSS/普通栈 & RW/NX & 状态可变但不可取指\\
heap & RW/NX & 动态对象默认不是代码\\
type 1 RAM direct-map & RW/NX & supervisor 访问物理页内容，未知洞不映射\\
零页与四个 guard & not-present & 把空指针和向下越界变成 \#PF\\
\bottomrule
\end{longtable}

NX 叶位只有在 \code{IA32\_EFER.NXE}=1 时才具有预期语义。内核先用 CPUID
扩展特性叶检查 NX，再设置 EFER.NXE；若硬件不支持就返回
\code{MemoryProtectionUnsupported}，不默默退化为全可执行。随后设置
\code{CR0.WP}，使 supervisor 写也遵守只读叶权限。

\section{TLB 是页表结果的处理器缓存}

处理器不会每次访问都重新走四级表。修改页表后，本核心可能继续使用旧 TLB
项；\code{invlpg} 使单页失效，写 CR3 或相应指令刷新更大范围。PCID 可以
减少地址空间切换的刷新，但需要额外代际管理。

多核系统每个核心有自己的 TLB。取消映射后必须向曾运行该地址空间的核心发送
shootdown，并等待确认，才能释放和复用物理页。否则另一个核心可能通过旧翻译
写入已经分配给新对象的页面。

\section{写保护故障证明权限不是软件自我声明}

页表查询函数能解码 RW=0，却不能证明处理器当前真的使用这张表，也不能证明
CR0.WP 已启用。v0.6 另外分配一帧，映射到
\code{0xFFFF800000100000}，权限为 present、supervisor、read-only、NX。
故障镜像在新 CR3 激活后执行 64 位写。

预期页故障错误码为 \code{0x3}：
\[
  \mathtt{0x3}=\underbrace{1}_{P}
                 +\underbrace{2}_{W/R},
\]
U/S、RSVD 和 I/D 都为零。P=1 把它与“忘记映射测试页”的错误码 2 区分；
W/R=1 把它与普通读取区分；CR2 必须精确等于高半区测试地址。只有随后进入同一
panic 并禁止 \code{READY}，才能说明权限执行、异常入口和停止边界全部连通。

\section{v1.1 用独立 KVA 层记录内核虚拟区间所有权}

\subsection{not-present 从来不等于 free}

早期内核的 heap、初始栈、IST、LAPIC 和写保护测试页都有固定地址。只要每种
资源恰好存在一份，链接脚本和常量表就能避免冲突；动态 Thread、可增长 slab、
临时设备缓冲和模块映射出现后，地址本身也成为必须申请、持有和归还的资源。

页表不能承担这份软件所有权。PTE 的 present 位只回答处理器现在能否把线性
地址翻译成物理地址。以下三种状态在硬件看来都可能是 present=0，软件含义却
完全不同：

\begin{longtable}{L{0.24\textwidth}L{0.22\textwidth}L{0.40\textwidth}}
\toprule
\textbf{软件状态} & \textbf{叶 PTE} & \textbf{能否交给新调用者}\\
\midrule
\endfirsthead
\toprule
\textbf{软件状态} & \textbf{叶 PTE} & \textbf{能否交给新调用者}\\
\midrule
\endhead
空闲地址 & not-present & 可以，由 KVA 分配器提交所有权\\
事务准备中 & not-present & 不可以，调用者正在取得物理后备或构造页表\\
guard page & 永久 not-present & 不可以，它属于栈区间并负责捕获越界\\
\bottomrule
\end{longtable}

若通过遍历 PTE 寻找 not-present 空洞，两个并发或交错事务都可能在首个映射
提交前得到同一地址；guard 也会被错误复用。KVA（kernel virtual address）
分配器因此只回答一个问题：“哪段虚拟页区间已经许诺给谁？”它不申请物理页、
不设置 present，也不选择 RW/NX 权限。

这种分层不是本项目发明的偶然术语。早期 UNIX 内核主要依赖静态映射；随着
内核模块、网络缓冲和动态对象增加，内核需要把不连续物理页组织成连续虚拟
区间。现代系统常把这类能力归入 vmalloc 一族，并加入树索引、延迟 TLB 刷新、
调试 guard 与并发回收。本项目先冻结串行、有界、可完整证明的区间状态机，
不复制尚无调用场景的性能层。

\subsection{32 TiB 窗口怎样落在四级页表中}

四级分页使用 48 位 canonical 线性地址：bit 63..48 必须全部复制 bit 47。
高半区因此从 \code{0xFFFF800000000000} 延伸到
\code{0xFFFFFFFFFFFFFFFF}。当前内核把 direct-map 与 KVA 排列为：

\begin{longtable}{L{0.20\textwidth}L{0.24\textwidth}L{0.24\textwidth}L{0.18\textwidth}}
\toprule
\textbf{区域} & \textbf{起点} & \textbf{末端（不含）} & \textbf{PML4 槽}\\
\midrule
\endfirsthead
\toprule
\textbf{区域} & \textbf{起点} & \textbf{末端（不含）} & \textbf{PML4 槽}\\
\midrule
\endhead
direct-map &
\code{FFFF888000000000} &
\code{FFFFC88000000000} &
273--400\\
KVA &
\code{FFFFC90000000000} &
\code{FFFFE90000000000} &
402--465\\
\bottomrule
\end{longtable}

一个 PML4 槽覆盖 512 GiB，64 个槽正好覆盖 32 TiB。KVA 起点的四级索引是
\((402,0,0,0)\)，exclusive 末端是 \((466,0,0,0)\)。起点、长度和末端都在
提交前按 64 位无符号数验证：页数乘 4096 不得溢出，末端加法不得回绕，末页
仍必须满足 48 位 canonical 规则。未来若启用五级分页 LA57，需要新的布局
迁移，不能让同一个 ABI 常量在两种符号扩展规则下静默改变含义。

窗口包含
\[
 \frac{32\ \mathrm{TiB}}{4\ \mathrm{KiB}}
 =8589934592=2^{33}
\]
个虚拟页。首个 4 KiB 页永久标记为 reservation。它不占物理页，也没有叶
映射；它只让窗口零偏移误用具有稳定的软件诊断，而不是直接成为第一个正常
对象。

\subsection{为什么 32 TiB 没有配一张逐页位图}

若每个虚拟页只用一位表示“占用”，元数据也需要
\[
 \frac{2^{33}}{8}=2^{30}\ \mathrm{bytes}=1\ \mathrm{GiB}.
\]
要区分 reservation 与 allocation 至少还要增加状态位。启动阶段实际同时
存在的 KVA 对象远少于数十亿页，为稀疏可能空间支付 GiB 常驻元数据没有比例
意义。

当前实现由调用者提供 256 项
\code{KernelVirtualAddressRangeDescriptor} BSS 数组。每项恰有三个 64 位
字段：

\begin{center}
\begin{tabular}{lll}
\toprule
字段 & 含义 & 有效约束\\
\midrule
\code{begin\_address} & 区间首虚拟地址 & 4 KiB 对齐、位于窗口内\\
\code{page\_count} & 区间页数 & 非零、换算字节不溢出\\
\code{kind} & Unused/Allocation/Reservation & 活动前缀不能为 Unused\\
\bottomrule
\end{tabular}
\end{center}

总大小为 \(256\times24=6144\) 字节。数组前部是按起始地址严格递增的活动
描述符，后部必须保持零地址、零页数和 Unused。256 限制的是“同时存在多少
区间”，不限制单个区间能有多少页。描述符耗尽返回
\code{MetadataExhausted}；窗口没有足够连续页返回
\code{OutOfVirtualAddressSpace}。区分这两种失败，才能判断应扩元数据还是
回收地址。

空闲区间不保存节点，而由窗口边界和相邻有主区间的缝隙推导：

\begin{center}
\begin{tikzpicture}[os diagram,node distance=0mm]
  \node[os state,minimum width=2.2cm] (r) {reservation};
  \node[os block,right=of r,minimum width=2.5cm] (f1) {free gap};
  \node[os state,right=of f1,minimum width=2.2cm] (a) {allocation};
  \node[os block,right=of a,minimum width=2.5cm] (f2) {free gap};
\end{tikzpicture}
\end{center}

删除一个 allocation 描述符后，两侧 gap 在模型上立即连成一个空闲区间，
不存在另一条空闲链需要合并。代价是申请要扫描至多 256 项，插入和删除要移动
尾部，复杂度为 \(O(N)\)。当前串行规模下，完整可检查的不变量比尚未被测量
需要的红黑树更重要。

\subsection{best-fit 必须对齐绝对虚拟页号}

请求页数 \(n\)、对齐页数 \(a\) 时，\(a\) 必须是二的幂。对一个字节起点为
\(g\) 的 free gap，先得到绝对虚拟页号
\[
 p=\frac{g}{4096},
 \qquad
 p_{\mathrm{aligned}}=(p+a-1)\mathbin{\&}\sim(a-1),
\]
再把它安全乘回字节地址。这里不能对
\((g-\mathrm{window\_begin})/4096\) 做相对对齐；“8 页对齐”的契约是返回
地址本身能被 32 KiB 整除，不是它与窗口起点的距离恰好整除。

每个 gap 都先检查对齐前缀、剩余页数和溢出。所有可容纳候选中选择原始 gap
页数最小者；同样大小时保留扫描得到的低地址。这是 best-fit：优先消耗最小
可用洞，减少把大洞用于小请求的概率。它不是消除碎片的数学保证，但对固定
上限提供确定、可由独立模型复现的结果。

候选地址、插入位置和计数增量全部验证后，实现才移动描述符并填写调用者输出。
任何失败都保持数组、统计和输出 range 原值。这样上层跨资源事务可以根据
明确状态执行回滚，不会拿到“返回失败但部分有效”的地址。

\subsection{reservation 与精确释放阻止所有权猜测}

\code{ReserveRange} 用来表达固定软件占用。它与普通申请共享页对齐、窗口、
canonical、溢出和重叠检查，但不能由 \code{TryRelease} 释放。重叠检查先于
元数据容量检查，因此即使 256 项已满，重复保留已有地址也仍报告
\code{RangeOverlap}，而不是被较表面的容量错误遮蔽。

释放必须同时匹配原描述符的起始地址、页数和 Allocation 类型。区间内部地址
返回 \code{AllocationNotFound}，正确起点但错误页数返回
\code{AllocationSizeMismatch}，reservation 返回 \code{ReservedRange}，
已经删除的地址再次释放仍返回 \code{AllocationNotFound}。实现不会向前寻找
“看起来像块头”的数据，也不会猜测调用者希望释放整个包含区间。

成功删除后，数组尾部向前压紧，最后一项恢复全零。累计统计始终满足
\[
 \mathrm{successful\ allocations}
 -
 \mathrm{releases}
 =
 \mathrm{active\ allocations}.
\]
当前分配页与保留页之和不得超过 \(2^{33}\)，当前值不得超过各自峰值。

\subsection{Validate 从唯一有主集合重建全部事实}

\code{Validate} 不申请临时内存。它先验证窗口本身，再沿活动前缀逐项证明：

\begin{enumerate}[label=\arabic*.]
  \item 类型只能是 Allocation 或 Reservation，页数非零，区间完整位于窗口；
  \item 当前起点不小于前一区间末端，因此允许紧贴但禁止倒序和重叠；
  \item 重新累计分配页、保留页、活动申请和 reservation 数；
  \item 活动前缀之后每个描述符都严格为 Unused、零地址、零页数；
  \item 从窗口首尾与相邻描述符重新计算 free page 和最大连续 gap；
  \item 把重建结果与缓存统计、累计申请/释放和峰值逐一比较。
\end{enumerate}

空闲 gap 没有第二份持久索引，所以不会出现“活动树正确、空闲链漏了一块”的
双结构漂移。若未来为性能加入平衡树或分段索引，它必须被视为可重建的派生
结构，完整验证仍应回到唯一所有权集合。

\subsection{跨三层创建必须正向提交、逆向撤销}

取得 KVA 只是跨层事务的第一步。一个带 guard 的动态对象遵循：

\begin{verbatim}
allocate KVA range
  -> allocate physical block or pages
  -> map data pages with explicit permissions
  -> publish object
\end{verbatim}

任何中间失败都要从已提交的最后一步开始撤销。正常销毁同样固定为：

\begin{verbatim}
unpublish object
  -> clear leaf PTEs and invalidate TLB entries
  -> return physical pages to buddy
  -> release exact KVA range
\end{verbatim}

顺序不能交换。若先归还物理页，旧 TLB 翻译仍可能写入已经交给新对象的页；
若先释放 KVA，另一个调用者可以在旧映射尚存时取得同一虚拟地址。guard page
从申请到释放始终没有叶映射，却与数据页一起属于同一个 KVA 区间，这正是
“not-present 不等于 free”的执行例子。

当前 \code{UnmapPage} 不只清除叶项。它先证明整条路径有效且属于当前页表根，
再按 PT、PD、PDPT 从子到父回收已经变空的独占分支。KVA 自检先用一个虚拟页
和一个 order-0 物理页暖机首条页表路径：撤销时回收 PT 与 PD，却保留可能仍
被进程 PML4 项引用的共享 PDPT。主事务随后复用该共享入口，重新创建并在最后
一张叶撤销时回收 PT 与 PD。因此物理页、buddy、KVA 与页表自身都进入可解释
的所有权基线；唯一保留的 PDPT 是共享协议，不是被忽略的泄漏。

\subsection{六页整机事务把 guard、权限和回收连成一条证据链}

预热后，QEMU 中的正式 KVA 自检执行：

\begin{enumerate}[label=\arabic*.]
  \item 申请 6 页、8 页对齐区间，得到
        \code{0xFFFFC90000008000}；
  \item 从 buddy 申请 order 2，即 4 个连续物理页；
  \item 相对页 0 和 5 保持 not-present，作为两个 guard；
  \item 把相对页 1--4 映射为 supervisor、RW、NX；
  \item 用 \code{QueryPage} 核对四个物理地址连续、页粒度和权限精确；
  \item 经首尾数据虚拟页写入两个 64 位模式并读回；
  \item 逆序撤销四个映射，归还整个 order-2 块，再精确释放六页 KVA；
  \item 比较三层基线，运行 buddy 与 KVA 的完整校验，并核对两段事务累计
        回收两张 PT、两张 PD、零张 PDPT，精确保留一张共享 PDPT。
\end{enumerate}

只有事务全部提交，冷路径才输出一次稳定快照：

\begin{verbatim}
[OS][KERNEL] KVA_ALLOCATOR_READY
[OS][KERNEL] KVA_WINDOW_BASE=0xFFFFC90000000000
[OS][KERNEL] KVA_WINDOW_SIZE_BYTES=0x0000200000000000
[OS][KERNEL] KVA_DESCRIPTOR_CAPACITY=0x0000000000000400
[OS][KERNEL] KVA_ACTIVE_DESCRIPTORS=0x0000000000000001
[OS][KERNEL] KVA_ALLOCATED_PAGES=0x0000000000000000
[OS][KERNEL] KVA_RESERVED_PAGES=0x0000000000000001
[OS][KERNEL] KVA_SUCCESSFUL_ALLOCATIONS=0x0000000000000002
[OS][KERNEL] KVA_RELEASES=0x0000000000000002
[OS][KERNEL] KVA_PEAK_ALLOCATED_PAGES=0x0000000000000006
[OS][KERNEL] KVA_SELF_TEST_MAPPED_PAGES=0x0000000000000004
[OS][KERNEL] KVA_SELF_TEST_GUARD_PAGES=0x0000000000000002
[OS][KERNEL] KVA_SELF_TEST_PASSED
[OS][KERNEL] PAGE_TABLE_RECLAIMED_LEVEL1_TABLES=0x0000000000000002
[OS][KERNEL] PAGE_TABLE_RECLAIMED_LEVEL2_TABLES=0x0000000000000002
[OS][KERNEL] PAGE_TABLE_RECLAIMED_LEVEL3_TABLES=0x0000000000000000
[OS][KERNEL] PAGE_TABLE_RETAINED_SHARED_LEVEL3_TABLES=0x0000000000000001
[OS][KERNEL] PAGE_TABLE_RECLAIM_SELF_TEST_PASSED
\end{verbatim}

申请、描述符移动、映射和释放热路径不逐步打串口。日志只证明最终事务和守恒，
不会让 115200 波特 I/O 改变分配器时序。

\subsection{三类宿主模型与真实处理器路径各自证明什么}

单元测试覆盖未初始化、无存储、零容量、canonical 空洞、保留重叠、best-fit、
绝对对齐、输出失败原子性、描述符与地址分别耗尽、精确释放错误和人为破坏
元数据。它证明边界和每个状态码，但不宣称页表真正执行。

集成测试串联真实 \code{PhysicalFrameAllocator}、
\code{KernelVirtualAddressAllocator} 与 \code{PageTableManager}，在宿主
模拟物理内存上执行同一六页生命周期。它核对 guard、四个物理身份、RW/NX、
写入模式和最终三层统计，证明模块组合没有遗漏回滚。

固定种子 \code{0x4B564152414E444F} 在 512 页独立逐页模型上执行 100000 步
申请与释放。参考模型自己扫描 gap、计算 best-fit 与对齐，不调用被测分配器；
每 257 步比较活动集合、累计统计、最大空洞并运行完整校验，结束时再随机顺序
排空所有申请。

QEMU 最后让真实 x86-64 MMU、TLB 和 load/store 执行同一契约。64 MiB 证明
低资源兼容路径，64 GiB 证明现代规格没有改变地址层语义。三类宿主测试加一条
处理器执行路径不是重复劳动：它们分别定位算法、组合、长序列不变量和硬件
效果。

当前 v1.2 KVA 边界仍然明确：最多 1024 个有主区间、操作为 \(O(N)\)、
没有内部锁。v1.1 的初始上界为 256；本次扩容用于同时容纳 512 个动态
Thread 栈及其他启动事务，不改变区间算法。
空中间页表回收已经由独立的页表所有权协议闭环；动态内核栈则成为第一个长期
持有 KVA、物理帧和叶映射的真实调度资源。性能树、per-CPU KVA cache 和延迟
shootdown 必须等并发契约与测量证据出现后再加入。

\section{v1.1 页表空分支回收：硬件树之上的软件所有权}

\subsection{四级结构来自地址空间扩张，不来自对象生命周期}

80386 的 32 位分页已经采用目录加页表的层次结构，避免为整个线性地址空间
预先保存逐页表项。PAE 把表项扩为 64 位并增加层级，x86-64 又让常见四级
模式用四组 9 位索引覆盖 48 位 canonical 线性地址。这个历史方向始终在解决
“怎样让硬件快速找到翻译”，没有替操作系统回答“最后一个映射消失后谁可以
释放中间表”。

在四级 4 KiB 模式中，每张表恰有 512 个 8 字节项，所以正好占一个 4096
字节物理页。各层覆盖范围为：

\begin{longtable}{@{}lll@{}}
\toprule
层 & 一个表项覆盖 & 一张完整表覆盖\\
\midrule
PT / level 1 & \(4\ \mathrm{KiB}\) & \(2\ \mathrm{MiB}\)\\
PD / level 2 & \(2\ \mathrm{MiB}\) & \(1\ \mathrm{GiB}\)\\
PDPT / level 3 & \(1\ \mathrm{GiB}\) & \(512\ \mathrm{GiB}\)\\
PML4 / level 4 & \(512\ \mathrm{GiB}\) & \(256\ \mathrm{TiB}\) canonical envelope\\
\bottomrule
\end{longtable}

CR3 给出根表物理地址。处理器从虚拟地址取出 \(47{:}39\)、\(38{:}30\)、
\(29{:}21\)、\(20{:}12\) 四组索引，最后把 \(11{:}0\) 当页内偏移。每个
非叶项主要告诉硬件下一张表的物理页和 P/RW/U/S 等遍历权限；叶项再决定
最终页的 NX、缓存、脏位和访问位等属性。这里没有父指针、引用计数、
\code{owner} 或“这张表由哪个地址空间申请”的字段。

这一缺失不是 x86 设计漏洞。硬件只需要在每次 load/store 时完成确定遍历；
对象生命周期属于操作系统策略。代价是软件不能从某个非叶项孤立地推出子表
是否独占。

\subsection{复制 PML4 项共享的是子树，不是子树副本}

当前进程地址空间为了共享内核映射，会把内核根中的高半 PML4 项按值复制到
新的进程 PML4。复制的 64 位值仍包含同一 PDPT 物理地址：

\begin{verbatim}
kernel PML4E ----+
                 +----> shared PDPT ---> PD ---> PT
process A PML4E -+
process B PML4E -+
\end{verbatim}

因此“内核根里这条分支已经没有叶映射”和“没有任何根引用该 PDPT”是两个
不同命题。若内核只清自己的 PML4E 并归还 PDPT，进程 A/B 的 PML4E 仍指向
原物理帧。buddy 可以马上把该帧交给堆或文件缓存；下一次在 A 中发生硬件
page walk 时，处理器会把新对象字节解释为页表项。这是页表帧
use-after-free，通常比普通悬空指针更难定位，因为损坏在另一个 CR3 被调度后
才显现。

相反，PDPT 内某个 PDPTE 位于真正共享的物理页中。清除它会被所有引用该
PDPT 的根共同观察到，所以它指向的空 PD 以及更低的 PT 可以安全释放。这个
差异给出本阶段的核心原则：

\begin{quote}
空只说明没有有效内容；回收还必须证明当前管理器拥有最后一个硬件引用。
\end{quote}

\subsection{三种根类型把隐含关系变成构造契约}

\code{PageTableManager} 构造时必须显式接收
\code{PageTableRootKind}，不能从虚拟地址、当前 CR3 或调用栈猜测：

\begin{longtable}{@{}p{0.18\textwidth}p{0.42\textwidth}p{0.30\textwidth}@{}}
\toprule
根类型 & 允许修改的分支 & 最深回收边界\\
\midrule
\code{Exclusive}
& 整棵树都只属于该管理器
& PT、PD、PDPT\\
\code{KernelShared}
& 内核根中的映射；下级父项变化对所有借用者可见
& PT、PD；PDPT 保持稳定\\
\code{Process}
& PML4[0] 下的用户程序分支和 PML4[255] 下的用户栈分支
& 程序低端克隆 PDPT 留到进程销毁；用户栈可回收到 PDPT\\
\bottomrule
\end{longtable}

进程低地址还有一个细节。PML4[0] 指向克隆的 PDPT，其中只有约定的
PDPT[1] 用户程序分支归进程修改；其余低端内容可能来自启动兼容映射。用户栈
放在 PML4[255] 的独占分支。对复制来的内核高半分支或其他借用分支调用
\code{MapPage}/\code{UnmapPage}，必须在申请帧和写表项之前返回
\code{SharedBranchMutationDenied}。

把根类型设为构造必选参数还有工程价值：新增调用点若没有回答所有权问题就
无法编译。相比默认成“最宽松”或根据地址猜测，这会把架构判断前移到代码
评审和单元测试。

\subsection{“空”必须表示其余 511 个原始项逐字为零}

一个表是否可回收，不能只数 Present 位。现代操作系统经常在 non-present
PTE 中保存 swap 类型、文件页状态或诊断标记；处理器不继续遍历，不代表软件
没有保存事实。当前项目尚未实现 swap，也仍采用更强定义：排除马上要解除的
父子边后，其余每个 64 位项必须等于零。

对普通 4 KiB 叶，判定链为：

\begin{verbatim}
PT except target PTE is all-zero
  -> PD except target PDE is all-zero
    -> root policy permits PDPT reclamation
      and PDPT except target PDPTE is all-zero
\end{verbatim}

一旦某层不空，更高层必然仍有有效子树，检查即可停止。2 MiB 页把 PDE 本身
作为叶并设置 PS；其数据后备通常对应 order 9 连续页。当前
\code{UnmapPage} 有意拒绝这种情况并返回
\code{UnexpectedLargePage}，不把不同粒度的数据所有权偷偷塞进 4 KiB
接口。

\subsection{撤销先完整预检，再一次性断开}

直接“边走边清”会产生难以回滚的半状态。假设已经清掉 PTE，走到 PD 时才
发现父项指向未拥有或越界物理帧；调用方既失去映射，又无法证明中间树安全。
当前撤销因此分为两阶段。

只读预检逐层完成：

\begin{enumerate}[label=\arabic*.]
  \item 检查虚拟地址 canonical 且按 4 KiB 对齐；
  \item 检查每个子表物理地址按页对齐，完整落在允许访问范围；
  \item 拒绝子表等于根或任一祖先的明显环；
  \item 用 buddy 的精确查询证明每张子表仍是活动 order-0 allocation；
  \item 保存四级父项地址和三个子表物理身份；
  \item 按原始 64 位值计算 PT、PD、PDPT 是否会在本次撤销后变空；
  \item 在修改任何表项前，再次预检所有计划释放的帧。
\end{enumerate}

预检失败返回 \code{InvalidTableFrame} 或
\code{TableFrameNotOwned}，叶项、调用方结果对象、TLB 与分配器统计全部
保持原样。预检成功后的提交顺序是：

\begin{verbatim}
clear PTE
clear PDE/PDPTE/PML4E that lead only to reclaimable empty tables
INVLPG target virtual address exactly once
zero and release PT, then PD, then permitted PDPT
\end{verbatim}

父项先断开，确保新的硬件遍历无法进入即将归还的子表；页帧再从子到父释放，
与创建依赖顺序相反。一次目标叶修改只需一次
\code{INVLPG}，为每个已释放中间表重复失效同一线性地址没有额外语义。
当前是单 BSP，所以还没有向其他 CPU 发送 TLB shootdown。

\code{PageTableUnmapResult} 分别返回 level 1、level 2、level 3 和总回收
表帧数。数据页由上层对象拥有，不由撤销接口释放，也不计入结果。这使调用者
可以精确区分“PTE 已清但共享表仍存在”和“本次真正回收了页表帧”。

\subsection{映射也必须能撤销父项权限与新表}

映射一张原先不存在的叶页，可能依次成功创建 PDPT、PD，却在申请 PT 时耗尽。
旧式实现若立即链接每张新表，只返回
\code{FrameAllocationFailed}，就会永久留下两个空页。另一个更隐蔽的例子是
用户映射：走过既有父项时需要把 U/S 提升为 1；若最后发现叶已经映射，临时
提升也不能残留。

当前每次非叶穿越记录一个 \code{TableMutation}：

\begin{verbatim}
entry                   address of the parent entry
original_entry          complete 64-bit value before mutation
table_physical_address  child table frame
table_created           whether this transaction allocated it
\end{verbatim}

只有叶项写入成功才提交事务。此前遇到帧耗尽、大页冲突、重复映射、损坏地址
或所有权错误时，先反序恢复每个父项完整原值，再反序清零并释放本事务创建的
表帧。保存完整值而非只保存 U/S 位，是为了不假设未来不会加入新的软件位。
若回滚自身遇到帧所有权或释放故障，返回
\code{RollbackFailed}，不能继续用最初错误掩盖已经不可证明的树状态。

这种局部日志与数据库 write-ahead log 的目的相似：事务先保存足以恢复的
旧事实，再改变共享结构。区别是它使用固定长度栈上记录，不分配内存，也没有
持久化恢复需求。

\subsection{查询和进程整体销毁也必须服从同一边界}

\code{QueryPage} 是只读接口，却不能无条件解引用页表地址。损坏或已释放的
父项同样会让宿主测试越界、让目标内核读取任意 direct-map 内容。因此查询在
每一级复用范围、对齐、环和精确 order-0 所有权验证。命中 4 KiB 页时返回
\[
  \mathrm{physical\ page\ base} + (\mathrm{virtual\ address}\bmod4096),
\]
不能像旧实现那样遗漏页内偏移；2 MiB 叶使用对应的大页偏移。

逐页撤销是正常细粒度路径，\code{ReleaseProcessRoot} 仍作为地址空间最终
销毁的兜底。它只遍历用户程序与用户栈的自有树，不进入借用的内核高半子树；
已被逐页撤销清空的分支自然跳过，尚存的用户叶和表由递归路径归还。每释放
一张子表都立即清父项，避免部分失败留下指向已归还帧的硬件指针。递归下降
同时携带从根到当前表的完整祖先物理地址链；若下级项回指根或任一祖先，
必须在释放任何相关表帧前返回 \code{InvalidTableFrame}。测试随后恢复原项并
重试销毁，用同一个对象证明拒绝路径没有偷改所有权状态。

\subsection{从单叶到十万步模型的证据梯度}

页表回收新增三类宿主测试。单元测试分别建立三种根，覆盖重复初始化、单叶
完整级联、同 PT 相邻叶、共享 PDPT 保留、进程借用拒绝、程序与栈分支差异、
申请耗尽回滚、所有权破坏与原址修复、越界表地址、查询祖先环、递归销毁
祖先回指和失败输出保持。

集成测试不在每轮重建世界，而是在同一分配器上连续执行 128 次共享根生命周期
和 64 次进程根创建、映射、逐页撤销与销毁。这样一个每轮泄漏一页的错误无法
被 fixture 清零隐藏；最后还让进程递归销毁一棵未逐叶排空的树。

固定种子 \code{0x5047545245434C4D} 在 1024 个虚拟页上执行 100000 步。
地址集合跨 4 个 PML4 项、每项 2 个 PDPT 项、多个 PD 与 PT；独立模型只保存
活动叶集合，并由层级索引重算应存在的 PT/PD/PDPT 数，不调用被测回收逻辑。
每 257 步逐页查询、比较表数和 buddy 统计，结束时排空全部映射。

QEMU 最后用真实 CR3、page walk、TLB 和访存验证暖机与四叶 KVA 事务。稳定
摘要为：

\begin{verbatim}
PAGE_TABLE_RECLAIMED_LEVEL1_TABLES = 2
PAGE_TABLE_RECLAIMED_LEVEL2_TABLES = 2
PAGE_TABLE_RECLAIMED_LEVEL3_TABLES = 0
PAGE_TABLE_RETAINED_SHARED_LEVEL3_TABLES = 1
\end{verbatim}

PT 与 PD 各两张来自暖机和主事务；PDPT 为零并精确保留一张，是共享根策略的
正面证据。映射、查询、撤销、空表扫描和随机测试都不逐项写串口，只在完整
初始化提交后输出一次摘要。至此 KVA 与动态栈不再靠暖机泄漏建立基线。

当前边界也必须写清：回收仅支持普通 4 KiB 叶；页帧元数据只证明精确
order-0 allocation，还没有 PageTable 类型标签；\code{KernelShared}
保守保留全部 PDPT，没有通用引用计数；没有 PCID、批量 unmap、远端
shootdown、页表锁或 RCU 读者。这些属于 Process/Thread 与并发地址空间模型
之后的独立问题，不应削弱本阶段已经闭合的单核所有权契约。

\section{v1.1 动态内核栈把地址许诺变成可回收调度资源}

\subsection{特权切换为什么需要一块不能自我销毁的栈}

x86 的任务状态段诞生于硬件任务切换时代。现代 64 位内核通常不使用 TSS
保存整份任务现场，但处理器从低特权级进入高特权级时，仍需要一个不受用户
RSP 控制的可信落脚点。当前系统在 CPL3 收到中断、异常或
\code{INT 0x80} 时，处理器从本 CPU 的 TSS.RSP0 取 Ring 0 栈顶，再压入
旧 SS、RSP、RFLAGS、CS 和 RIP。汇编入口继续在同一栈上保存通用寄存器。

这形成一个重要的生命周期悖论：进程终止代码正在使用的资源，恰好也是该进程
应当释放的资源。若退出处理函数直接清除当前栈的 PTE，下一条函数返回、局部
变量访问或寄存器恢复就会触发内核页故障；若先把物理页归还 buddy，旧 TLB
翻译还可能写入已经交给其他对象的帧。正确答案不是“销毁函数更小”，而是把
终止分成逻辑终止和物理回收两个时刻。

v0.9 为了先证明抢占调度，给四个 PCB 各编译进一块固定 20 KiB 存储，并在
低端留下一个 guard。这个实现没有地址申请、没有物理帧所有权，也没有真正的
销毁状态；容量和 PCB 槽位在链接时绑定。v1.1 保留四进程教学负载，却把栈
本身迁移成跨 KVA、buddy 和页表三层的动态对象。这样 v1.2 把 PCB 拆为
Process/Thread 时，可以迁移所有者，而不必再次发明栈资源协议。

\subsection{六页布局同时表达容量和两侧故障边界}

一个动态内核栈取得 6 个连续虚拟页，但只有中间 4 页存在叶映射：

\begin{center}
\begin{tikzpicture}[os diagram,node distance=0mm]
  \node[os hardware,minimum width=1.9cm] (lower) {lower guard\\4 KiB};
  \node[os state,right=of lower,minimum width=2.1cm] (d0) {data 0\\4 KiB};
  \node[os state,right=of d0,minimum width=2.1cm] (d1) {data 1\\4 KiB};
  \node[os state,right=of d1,minimum width=2.1cm] (d2) {data 2\\4 KiB};
  \node[os state,right=of d2,minimum width=2.1cm] (d3) {data 3\\4 KiB};
  \node[os hardware,right=of d3,minimum width=1.9cm] (upper) {upper guard\\4 KiB};
\end{tikzpicture}
\end{center}

设 KVA 返回的区间起点为 \(b\)，页大小 \(P=4096\)，则：
\[
\begin{aligned}
  \mathrm{lower\ guard} &= b,\\
  \mathrm{mapped\ begin} &= b+P,\\
  \mathrm{stack\ top} &= b+5P,\\
  \mathrm{upper\ guard} &= b+5P,\\
  \mathrm{range\ end} &= b+6P.
\end{aligned}
\]
可用容量为 \(4P=16384\) 字节，完整虚拟区间为 \(6P=24576\) 字节。x86-64
栈向低地址增长，所以 TSS.RSP0 指向 \(\mathrm{stack\ top}\)，处理器第一次
压栈落入 data 3。继续向下越过 data 0 会命中 lower guard。upper guard
不负责捕获正常的向下溢出；它负责隔开相邻对象，并让错误的正偏移、范围计算
或把“栈顶”误作可写首地址时立即得到 not-present 故障。

两个 guard 从申请到释放都属于同一 KVA 描述符，却始终没有 PTE。这正是
“not-present 不等于 free”的长期实例。若只保留 lower guard 而让下一对象
紧贴栈顶，错误的向上访问可能静默写进另一个内核对象；双 guard 用额外 4 KiB
虚拟地址换取对称且稳定的故障边界，并不消耗额外物理 RAM。

\subsection{虚拟连续不要求四个物理帧连续}

\code{KernelStack} 保存一个六页
\code{KernelVirtualAddressRange}、四个 \code{PhysicalFrame}、已映射
页数和 active 状态。管理器对四个 data 页分别调用 order-0 物理分配：
\[
  V_{b+P+iP}\longmapsto F_i,\qquad i\in\{0,1,2,3\}.
\]
\(F_i\) 可以位于任意四个可用 PFN；页表把它们组织成连续的 16 KiB 虚拟栈。
若为每个小栈申请 order 2 连续块，长期运行后即使总空闲页很多，也可能因外部
碎片找不到四页连续物理区。order 0 更符合栈只要求虚拟连续的真实契约。

每帧在映射前通过 direct-map 清零。清零既避免新进程观察到旧内核栈残留，也
让首次构造现场以确定状态开始。销毁时再次清零后才归还 buddy，使后续所有者
不会继承寄存器、地址或用户参数。叶权限固定为 supervisor、RW、NX、
cache enabled；任何 user、executable 或 cache-disabled 偏差都由完整校验
拒绝。

管理器有 512 个长期槽，容量与 1024 项 KVA 描述符容量显式分开。槽索引由
Thread 持有，不等于 PID、TID 或数组身份；栈管理器本身不理解调度状态或
用户地址空间。统计记录：

\begin{longtable}{L{0.31\textwidth}L{0.53\textwidth}}
\toprule
\textbf{字段} & \textbf{必须成立的含义}\\
\midrule
\code{active\_stack\_count} & 当前 active 槽数量\\
\code{active\_mapped\_page\_count} & active 数量乘 4\\
\code{active\_guard\_page\_count} & active 数量乘 2，由前两者派生\\
\code{successful\_creation\_count} & 从初始化以来完整发布的栈数量\\
\code{destruction\_count} & 完整撤销并清空槽的数量\\
\code{peak\_active\_stack\_count} & 历史同时活动栈峰值\\
\code{peak\_active\_mapped\_page\_count} & 历史物理后备页峰值\\
\bottomrule
\end{longtable}

累计计数必须满足：
\[
  \mathrm{creations}-\mathrm{destructions}
  =\mathrm{active\ stacks},
\]
\[
  \mathrm{active\ mapped\ pages}=4\times\mathrm{active\ stacks}.
\]
峰值只增不减，当前值不得超过峰值。64 位计数在提交前检查溢出，不能让一次
成功操作把守恒式绕回零。

\subsection{创建是一笔直到发布前都可完整撤销的事务}

\code{TryCreate(slot)} 先验证管理器、槽位、依赖分配器和统计，再执行：

\begin{enumerate}[label=\arabic*.]
  \item 向 KVA 申请 6 页、1 页对齐的完整区间；
  \item 查询 6 个虚拟页，证明没有任何遗留叶映射；
  \item 逐页申请四个 order-0 物理帧并验证 direct-map 可访问；
  \item 清零当前帧，以 supervisor、RW、NX 权限映射对应 data 页；
  \item 核对两个 guard 仍 not-present、四个叶映射的物理身份和权限；
  \item 把候选对象复制到目标槽，最后设置 active 并更新统计。
\end{enumerate}

候选对象在最后一步之前不对调度器可见。第 \(j\) 个物理帧申请或映射失败时，
\code{RollbackCreation} 按相反顺序清除已经建立的叶项、清零并归还已经取得
的帧，最后精确释放六页 KVA。伪代码为：

\begin{verbatim}
range = kva.allocate(6)
for each data page:
    frame = buddy.allocate(order0)
    zero(frame)
    map(data_va, frame, supervisor | rw | nx)
validate(candidate)
publish(slot, candidate)

on any failure:
    unmap(mapped pages in reverse order)
    zero and release acquired frames in reverse order
    kva.release(exact range)
\end{verbatim}

回滚本身也可能失败，因此接口区分原始的
\code{FrameAllocationFailed} 或 \code{PageMappingFailed}，以及回滚时的
\code{RollbackFailed}。后一种状态说明资源集合已经无法证明，调用方不能把
它当成普通“内存不足”继续运行。失败输出不发布半有效 \code{KernelStack}，
目标槽仍保持全零。

\subsection{精确 KVA 所有权阻止部分销毁}

仅验证六页都在 KVA 窗口内还不够。若其他错误代码提前释放了栈的 KVA
描述符，PTE 和物理帧可能仍然存在；此时销毁函数若先撤映射、归还物理帧，
最后才发现 KVA 不再归自己，就会留下无法重试的半销毁对象。

\code{KernelVirtualAddressAllocator::OwnsAllocation} 因而要求起点、页数和
Allocation 类型与活动描述符精确匹配。物理分配器的同名只读查询只接受精确
order-0 活动块；大块内部首帧或已释放帧都不满足。\code{ValidateStack} 在
任何破坏性步骤前同时验证两层所有权、两个 guard、四个叶映射、权限和四个
物理地址互异。只有预检完整通过，销毁才逆序：

\begin{verbatim}
unmap data pages from high to low
zero and release frames from high to low
release the exact six-page KVA allocation
clear the slot and commit statistics
\end{verbatim}

当前单 BSP、关中断提交边界让预检与撤销之间没有另一个管理器调用者。SMP
版本必须把“验证后执行”放入同一个锁或代际协议，否则预检成功不能阻止另一个
CPU 在提交前改变页表或 KVA 所有权。

\subsection{高半区共享让所有进程在同一虚拟地址看见栈}

动态栈位于 KVA 窗口，对应 PML4 槽 402--465。创建进程根时，页表管理器复制
永久内核 PML4 的高半区项目；因此每个进程 CR3 都通过共享的内核页表子树看到
相同栈映射。栈不是用户地址空间中的私有叶页，也不需要为四个进程重复建立四份
映射。

共享带来两条严格顺序：

\begin{itemize}
  \item 创建栈必须发生在克隆进程根之前，使新根取得已经存在的高半区路径；
  \item 销毁用户地址空间只能释放低半区独占子树，不能递归释放 KVA 的共享
        高半区表。
\end{itemize}

当前高半区页表项本身由所有进程共享，所以在永久内核 CR3 下增加或清除叶项，
其他进程根会观察同一结构。单核路径通过 CR3 切换和 \code{INVLPG} 处理当前
TLB；SMP 后必须跟踪哪些 CPU 运行过相关根，并在物理帧复用前完成 shootdown。

\subsection{176 字节首次现场必须完全落在最后一个 data 页}

\code{UserPrivilegeFrame} 长 176 字节。创建进程时，运行时读取管理器发布的
栈顶 \(t\)，按 16 字节 ABI 对齐后在其下方构造首次现场：
\[
  f=\operatorname{alignDown}(t-176,16).
\]
管理器的 \code{Contains(slot,f,176)} 证明半开区间
\([f,f+176)\) 完整落在四个 data 页中。当前布局使它落在最高 data 页，写入
物理内存时需使用该页对应的独立 frame 和页内偏移，不能假设四个 PFN 连续。

首次派发前，\code{ActivateProcess} 同时验证保存帧仍属于该槽、切换到进程
CR3，并把 TSS.RSP0 设置为该动态栈顶。抢占切换时返回的另一个帧地址和新的
RSP0 属于同一个目标进程；若只换 CR3 不换 RSP0，下次用户态陷入会把现场压到
旧进程的栈上。

\subsection{安全点回收把逻辑死亡和物理释放分开}

用户 exit 或用户异常发生时，C++ 处理函数正在终止进程的动态栈上运行。
\code{CompleteCurrentProcess} 只完成以下工作：

\begin{enumerate}[label=\arabic*.]
  \item 保存终止现场并把调度状态改为 Terminated；
  \item 关闭描述符并选择后继；
  \item 切换到永久内核 CR3，销毁进程独占的用户页和页表；
  \item 有后继时切换其 CR3/RSP0 并返回其保存帧；无后继时恢复默认 RSP0；
  \item 由汇编恢复路径最终返回永久启动栈。
\end{enumerate}

此时终止栈仍映射且仍由管理器持有。只有
\code{OsKernelEnterScheduledProcess} 返回到永久启动栈之后，
\code{ReapTerminatedKernelStacks} 才扫描 Terminated 槽。回收器先读取真实
RSP，并用 \code{Contains} 检查一个 64 位探针范围；只要当前 RSP 仍在目标
栈内，就拒绝销毁。通过检查后才调用 \code{TryDestroy}，成功后把 PCB 中的
保存帧指针清为 null。

\begin{center}
\begin{tikzpicture}[os diagram]
  \node[os state] (user) {Ring 3\\运行};
  \node[os block,right=of user] (exit) {exit/\#UD/\#PF\\逻辑终止};
  \node[os state,right=of exit] (switch) {永久 CR3\\用户空间回收};
  \node[os block,below=of switch] (boot) {汇编返回\\永久启动栈};
  \node[os state,left=of boot] (reap) {RSP 安全检查\\动态栈销毁};
  \draw[os arrow] (user) -- (exit);
  \draw[os arrow] (exit) -- (switch);
  \draw[os arrow] (switch) -- (boot);
  \draw[os arrow] (boot) -- (reap);
\end{tikzpicture}
\end{center}

“当前 CR3 已经不是进程根”不能替代 RSP 检查，因为动态栈属于共享高半区，在
永久根里同样存在。“进程状态已经 Terminated”也不能替代检查，因为状态是
调度事实，不是 CPU 当前栈指针的物理事实。安全点必须由真实处理器状态证明。

\subsection{完整校验从槽集合重建跨层事实}

\code{KernelStackManager::Validate} 不只比较计数。它遍历全部 512 槽并
重新证明：

\begin{enumerate}[label=\arabic*.]
  \item inactive 槽的 range、frame、页数和 active 位全部为零；
  \item active range 为 6 页、页对齐、无加法溢出且首尾均 canonical；
  \item KVA 分配器精确拥有这一个 allocation；
  \item lower/upper guard 的查询结果都是 NotMapped；
  \item 四个 data 页各自映射到记录的物理帧，页粒度为 4 KiB；
  \item 权限精确为 supervisor、RW、NX、普通缓存；
  \item 四帧均由 buddy 以 order 0 持有、direct-map 可访问且彼此不同；
  \item 重建的活动数、映射页数、累计差和峰值与缓存统计一致；
  \item 最后再运行 KVA 与 buddy 自身的完整不变量检查。
\end{enumerate}

跨层校验让一种资源的“局部合法”不能掩盖另一层漂移。例如一个 PTE 仍 present
但 KVA 描述符已丢失，在页表看来可以访问，在 buddy 看来帧也已分配，只有联合
所有权协议能判为损坏。

\subsection{三层测试和三条整机路径怎样分工}

单元测试覆盖初始化失败、非法槽、重复创建/销毁、地址窗口耗尽、物理帧耗尽、
页表冲突、创建回滚、读出、范围包含、双 guard、权限、非连续物理帧、统计和
人为破坏。测试先外部释放一个物理帧，证明 PTE 仍在也不能冒充所有权；原址
重新申请后再外部释放 KVA，二者都要求 \code{Validate} 与销毁在破坏性操作
前拒绝。最后修复 KVA、篡改并恢复 PTE，只有三层重新一致后才允许安全销毁。

集成测试创建四个栈，逐一在最高 data 页构造 176 字节现场，并建立新的进程
页表根。它从该根再次查询所有 guard 与 data 映射，证明高半区共享不是单一
永久 CR3 下的偶然现象。逆序销毁后，四个物理帧全部清零，buddy 空闲数和
KVA 页数恢复空闲基线；共享页表根只保留一张仍可能被进程 PML4 引用的 PDPT，
临时 PT/PD 均已回收。创建、销毁均为 4，峰值活动栈为 4，峰值映射页为 16。

固定种子 \code{0x4B535441434B524E} 在 64 个槽和 4096 页独立所有权模型上
执行 100000 次随机创建/销毁。操作类型取随机数高位，槽索引使用下一次独立
随机值，避免线性同余发生器低位相关性伪造覆盖。参考模型自己寻找六页
best-fit 空洞；每 257 步比较映射、统计和三层不变量，最后随机状态全部排空。

QEMU 正常路径、用户 \#UD 隔离路径和用户 \#PF 隔离路径都必须经过同一动态栈
创建与回收协议。正常四进程路径要求每个 lower guard、栈顶和 upper guard
地址各出现一次；故障路径要求异常进程也回到安全点并释放栈。整机冷路径输出：

\begin{verbatim}
[OS][KERNEL] KERNEL_STACK_MANAGER_READY
[OS][KERNEL] KERNEL_STACK_SLOT_CAPACITY=0x0000000000000200
[OS][KERNEL] KERNEL_STACK_MAPPED_PAGES=0x0000000000000004
[OS][KERNEL] KERNEL_STACK_GUARD_PAGES=0x0000000000000002
[OS][KERNEL] KERNEL_STACK_SIZE_BYTES=0x0000000000004000
...
[OS][KERNEL] KERNEL_STACK_ACTIVE_STACKS=0x0000000000000000
[OS][KERNEL] KERNEL_STACK_SUCCESSFUL_CREATIONS=0x...
[OS][KERNEL] KERNEL_STACK_DESTRUCTIONS=0x...
[OS][KERNEL] KERNEL_STACK_PEAK_ACTIVE_STACKS=0x...
[OS][KERNEL] KERNEL_STACK_PEAK_MAPPED_PAGES=0x...
[OS][KERNEL] KERNEL_STACK_RESOURCES_RECLAIMED
\end{verbatim}

这些日志在准备阶段、每 Thread 创建结果和最终聚合点输出，不在页申请、映射、
时钟中断或回收循环中逐步输出。宿主 QEMU runner 继续给每行添加单调时间；
来宾使用 PIT 单调 tick 表达自己的时间域。丰富日志在这里意味着关键状态可
关联、计数可守恒，而不是把高频事件冲进 115200 波特串口。v1.2 的累计值随
4/4、64/128、256/512 三档 Process/Thread 容量事务变化，因此不再伪装成
单一常数；验收要求 successful creations 与 destructions 相等、active 为零，
并要求 peak mapped pages 等于 peak active stacks 的四倍。

\subsection{当前完成边界与下一资源问题}

动态栈已经消除了编译期四份 20 KiB 进程栈存储。v1.1 最初由四个 PCB 保存
管理器槽索引，存储上限为 256；v1.2 已把所有者迁移为 Thread，并把槽上限
提升为 512。管理器仍是单 BSP、无内部锁；栈大小固定 16 KiB，不支持按需
增长，也没有每 CPU emergency stack 或 stack canary。空 PT/PD 与独占
PDPT 回收已经完成，并由页表帧统计证明动态栈反复创建/销毁不会永久增长
中间表基础设施。v1.1 的最后一步进一步把这些局部结论收束为通用引用计数、
可组合的作用域回滚与跨层资源快照。

\section{v1.1 通用资源生命周期把“已经释放”变成可验证结论}

\subsection{资源管理的历史困难不是缺少一个 free 函数}

一个内核对象很少只拥有一块同质内存。动态内核栈同时占有 KVA 区间、四个
物理页、四张叶映射和管理器槽；用户进程还会占有页表根、用户代码页、用户栈、
描述符和调度现场。创建函数执行到一半失败时，只能释放已经取得的部分，而且
必须按依赖关系的反方向撤销：

\begin{center}
\begin{tikzpicture}[os diagram]
  \node[os state] (kva) {六页 KVA\\软件所有权};
  \node[os state,right=of kva] (frames) {四个 frame\\物理所有权};
  \node[os state,right=of frames] (pte) {四个 PTE\\地址翻译};
  \node[os state,below=of frames] (slot) {KernelStack 槽\\对外发布};
  \draw[os arrow] (kva) -- (frames);
  \draw[os arrow] (frames) -- (pte);
  \draw[os arrow] (pte) -- (slot);
  \draw[os arrow,dashed] (slot) -- node[right]{失败时反向} (pte);
  \draw[os arrow,dashed] (pte) -- (frames);
  \draw[os arrow,dashed] (frames) -- (kva);
\end{tikzpicture}
\end{center}

若先归还 frame 而 PTE 仍 present，原虚拟地址会指向后来取得该 frame 的
新对象；若先释放 KVA，另一个调用者可能取得同一虚拟区间，而旧 PTE 尚未
撤销。四个步骤都调用过某种 release 仍不够，顺序本身就是安全属性。

早期 C 系统常用逐层 \code{goto fail\_map}/\code{fail\_frame} 标签表达
展开。它没有语言运行时依赖，也能精确工作，但创建步骤增加时必须同步修改多个
标签；漏掉一条清理或把标签跳错一级，通常只在稀有故障注入中出现。普通 hosted
C++ 通过 RAII、标准容器和异常展开减轻这类重复；freestanding Kernel 没有
理由引入异常运行时，更不能把页表、锁和设备请求都误当成普通堆内存。

因此本阶段不是寻找一个“自动释放所有东西”的黑箱，而是建立三种互补机制：

\begin{enumerate}[label=\arabic*.]
  \item \code{ScopeRollback} 记录当前事务已经完成的可逆步骤；
  \item \code{ReferenceCounter} 决定共享对象何时取得最后释放资格；
  \item \code{ResourceSnapshot} 在较长生命周期两端证明跨管理器当前量恢复。
\end{enumerate}

前者负责单次调用内的失败展开，中者负责多个所有者之间的销毁时机，后者负责
发现局部检查互相抵消或遗漏的泄漏。三者不能相互替代。

\subsection{ScopeRollback 是无分配的显式小事务}

\code{ScopeRollback} 不持有动态容器。调用者提供固定
\code{ScopeRollbackAction} 数组，每项只保存一个
\code{bool (*)(void*) noexcept} 操作和一个上下文指针。它的状态机为：

\begin{center}
\begin{tikzpicture}[os diagram]
  \node[os state] (u) {Uninitialized};
  \node[os state,right=of u] (a) {Armed};
  \node[os state,above right=of a] (c) {Committed};
  \node[os state,below right=of a] (r) {RolledBack};
  \draw[os arrow] (u) -- node[above]{Initialize} (a);
  \draw[os arrow] (a) edge[loop above] node{TryPush} (a);
  \draw[os arrow] (a) -- node[above left]{Commit} (c);
  \draw[os arrow] (a) -- node[below left]{TryRollback / 析构} (r);
\end{tikzpicture}
\end{center}

外部存储有两个重要原因。第一，低内存失败路径不能为了登记“稍后释放内存”
再次申请内存，否则诊断机制会重入原故障。第二，动作容量成为调用点可审计的
上界；当创建过程新增一步时，测试和代码审查能直接看到容量是否同步改变。

典型控制流可以抽象为：

\begin{verbatim}
ScopeRollbackAction actions[capacity];
ScopeRollback rollback;
rollback.Initialize(actions, capacity);

acquire A;
rollback.TryPush(ReleaseA, context_a);

acquire B;
rollback.TryPush(ReleaseB, context_b);

publish candidate;
rollback.Commit();
\end{verbatim}

未执行 \code{Commit} 时离开作用域，析构从最后一个动作向第一个动作执行。
显式 \code{TryRollback} 用于调用者需要检查清理结果的失败分支；析构是防止
遗漏的最后安全网。回滚动作之一返回 false 时，实现仍继续执行更早动作，
最后报告 \code{ActionFailed}。这是有意选择：一张 PTE 撤销失败不意味着应
故意保留所有其他可独立归还的 frame 和 KVA。上层收到失败后仍不能把候选对象
当作干净状态继续使用。

\subsection{九项动态栈补偿动作如何推导}

动态栈创建取得一次 KVA 和四组“物理页加叶映射”。在取得顺序中，每个数据页
先申请并清零 frame，再建立映射；对应的最大补偿动作数是：
\[
  1_{\mathrm{KVA}}+
  4\times
  \left(1_{\mathrm{frame}}+1_{\mathrm{PTE}}\right)=9.
\]

\begin{center}
\begin{tabular}{cll}
\toprule
登记次序 & 正向步骤 & 回滚时的动作\\
\midrule
0 & 取得六页 KVA & 精确释放 KVA allocation\\
1,3,5,7 & 取得一个 order-0 frame & 清零并释放该 frame\\
2,4,6,8 & 建立一个数据叶映射 & 撤销对应 PTE\\
\bottomrule
\end{tabular}
\end{center}

动作数组按正向取得顺序登记，执行时从 8 走到 0，所以每个已映射页先 unmap，
再归还其 frame，最终才释放 KVA。若第三个 frame 申请失败，数组中只存在
KVA、前两个 frame 和前两个 PTE 的动作；未取得资源没有虚构的补偿项。候选
\code{KernelStack} 仍是局部值，目标槽只有完整验证后才覆盖，因此回滚期间
其他代码看不到半对象。

销毁既有活动栈没有直接复用同一事务，因为它必须先证明所有权仍一致，再执行
不可回退的正式销毁。创建失败展开和已发布对象销毁是不同状态转换；把二者
强行塞进一个通用函数会模糊“候选从未可见”和“活动对象正在消失”的边界。

\subsection{ReferenceCounter 冻结强引用的最小语义}

共享对象需要回答“还有多少所有者保证其存活”。本阶段以明确的 64 位状态机
定义强引用：

\[
R\in[0,2^{64}-1],\qquad
\operatorname{Acquire}(R)=R+1\ \text{only if}\ 0<R<2^{64}-1.
\]

释放语义为：
\[
\operatorname{Release}(R)=
\begin{cases}
(R-1,\mathrm{false}), & R>1,\\
(0,\mathrm{true}), & R=1,\\
\mathrm{failure}, & R=0.
\end{cases}
\]

\code{Start(initial)} 只允许在未启动状态设置非零初值。\code{TryAcquire}
拒绝从零复活，因为最后释放者一旦观察到零，就可能已经开始销毁相关页和锁；
此后把计数加回一不能恢复对象内容。\code{TryAcquire} 也拒绝最大值上溢，
否则无符号回绕会把最多所有者误写成零。

\code{TryRelease(bool\& is\_last)} 只有成功时才修改输出。这个看似细小的
失败原子性十分关键：若计数已经为零而函数先把 \code{is\_last} 写成 true，
调用者可能在失败处理里再次销毁对象。计数和输出必须作为一个逻辑结果提交。

这里的 v1.1 实现有意不是 atomic。v1.2 仍是单 BSP，基础原语先冻结可测试的所有权
语义。多 CPU 下不能只把字段替换成 \code{atomic<uint64\_t>} 就宣告正确：
最后释放前的对象写入需要怎样发布，零后获取如何与销毁互斥，弱引用如何升级，
析构与 RCU 或锁如何配合，都必须一起定义。后来的 v1.4 动态 KernelObject
选择管理器锁串行化强引用和最后释放，并明确不声称已有弱引用或 RCU；真正的
多核内存序仍留给 v2 对应阶段。

\subsection{ResourceSnapshot 为什么排除历史累计量}

一次资源事务成功创建再销毁后，当前活动量应恢复，但历史计数合理增加。若把
\code{successful\_allocations} 或峰值放入泄漏比较，每一个正确完成的事务
都会被误报。因此快照只保存稳定当前所有权，字段分成六组：

\small
\noindent\begin{tabularx}{\textwidth}{L{0.14\textwidth}L{0.12\textwidth}X}
\toprule
组 & 字段数量 & 当前量\\
\midrule
frame & 4 & managed、free、allocated、reserved（unavailable 由差值推导）\\
buddy & 1 & active blocks\\
heap & 4 & capacity、consumed、active requested bytes、allocation count\\
KVA & 7 & capacity/free/allocated/reserved page、活动描述符/申请/保留\\
kernel stack & 4 & slot capacity、active stack、mapped page、guard page\\
未来对象 & 6 & Process、Thread、FileDescription、Vnode、CachePage、BlockRequest\\
\midrule
合计 & 26 & 使用低 26 位差异掩码\\
\bottomrule
\end{tabularx}
\normalsize

未来六项在 v1.1 为零，但字段位置现在就固定，使后续阶段能接入同一个比较器，
不必另造一条“进程计数看起来归零”的旁路验收。快照不是永久用户 ABI；字段
语义改变仍需 ADR 和测试同步更新。

\subsection{内部守恒先于前后比较}

两个同样损坏的快照可能逐字段相等。若只比较 before 和 after，进入事务前
已经存在的错误会被当作“没有变化”。因此构造和比较都先检查每个管理器的
内部方程：

\[
\begin{aligned}
F_{\mathrm{accounted}}
  &=F_{\mathrm{free}}+F_{\mathrm{allocated}}
    +F_{\mathrm{reserved}}
    \le F_{\mathrm{managed}},\\
F_{\mathrm{unavailable}}
  &=F_{\mathrm{managed}}-F_{\mathrm{accounted}},\\
K_{\mathrm{capacity}}
  &=K_{\mathrm{free}}+K_{\mathrm{allocated}}+K_{\mathrm{reserved}},\\
H_{\mathrm{consumed}} &\le H_{\mathrm{capacity}},\\
S_{\mathrm{mapped}} &=4S_{\mathrm{active}},\\
S_{\mathrm{guard}} &=2S_{\mathrm{active}},\\
S_{\mathrm{KVA}} &=6S_{\mathrm{active}}.
\end{aligned}
\]

这里的 managed 是从零到最高受管物理地址的页跨度，不等于全部可分配 RAM。
E820 保留洞、MMIO 窗口与不足整页的边界都保持 unavailable。快照用
managed 与三类当前所有权的差值推导 unavailable，既不浪费差异位，也仍能
检测任何跨状态误转移。这个区别在 64 GiB QEMU 的 3--4 GiB PCI 洞上会直接
影响验收；把不等式误写成等式会错误拒绝合法容量布局。

buddy 活动页和活动块不能超过 frame allocated；KVA 活动对象数不能超过
描述符容量。所有加法、乘法和类型转换先做 64 位溢出检查。快照构造先写局部
候选，完整验证成功后才覆盖调用者输出；比较也先产生局部差异，失败不交付
半份掩码。

差异位 \(i\) 对应字段 \(i\)：
\[
  D=\bigvee_{i=0}^{25}
  \left([B_i\ne A_i]\ll i\right),\qquad
  C=\operatorname{popcount}(D).
\]
当前不依赖标准库 \code{popcount}，而是显式遍历 26 个字段；字段总数由
\code{static\_assert} 保证小于 64，避免未来新增字段后左移越界。零差异必须
同时满足 \(D=0\) 和 \(C=0\)。

\subsection{目标机有两道不同的生命周期边界}

第一道位于内存初始化尾部。管理器在保留的最后一个槽创建真实六页动态栈，
外层 \code{ScopeRollback} 登记销毁动作；随后验证栈布局、引用计数状态机和
回滚状态，显式执行回滚，再拍摄 after 快照。只有 26 字段恢复才输出：

\begin{verbatim}
[OS][KERNEL] RESOURCE_LIFECYCLE_READY
[OS][KERNEL] RESOURCE_SNAPSHOT_TRACKED_FIELDS=0x000000000000001A
[OS][KERNEL] RESOURCE_SNAPSHOT_CHANGED_FIELDS=0x0000000000000000
[OS][KERNEL] REFERENCE_COUNTER_SELF_TEST_PASSED
[OS][KERNEL] SCOPE_ROLLBACK_SELF_TEST_PASSED
[OS][KERNEL] RESOURCE_SNAPSHOT_SELF_TEST_PASSED
\end{verbatim}

第二道跨越真实四进程运行。创建第一个进程前拍摄 before；Shell、管道、
文件系统、抢占和用户异常隔离全部完成，汇编返回永久启动栈并回收所有终止栈
之后拍摄 after。任一字段变化使 \code{ProcessRuntime} 返回
\code{ResourceLeakDetected}，禁止输出：

\begin{verbatim}
[OS][KERNEL] PROCESS_RESOURCES_RECLAIMED
[OS][KERNEL] RESOURCE_SNAPSHOT_PROCESS_LIFECYCLE_PASSED
\end{verbatim}

启动期自检增加一次栈创建和销毁，四进程再增加四次，所以正常整机累计值为
五次创建、五次销毁；自检栈在进程阶段前已经释放，峰值仍是四个活动栈和
十六个映射页。区分“累计五次”和“峰值四个”正是当前量与历史量分离的实例。

\subsection{四层测试怎样避免同源自证}

\noindent\begin{tabularx}{\textwidth}{L{0.16\textwidth}L{0.18\textwidth}X}
\toprule
层级 & 输入规模 & 主要证明\\
\midrule
原语单元 & 边界与故障枚举 & 状态、输出原子性、逆序和幂等拒绝\\
固定种子随机 & 100000 个事务 & 与独立动作/资源模型逐步一致\\
组合集成 & 四个真实动态栈 & frame/buddy/KVA/PTE/stack 同时变化并恢复\\
QEMU 系统 & 256 MiB 与 64 GiB & 真实 CPU、页表、Shell、IPC、文件系统与回收\\
\bottomrule
\end{tabularx}

固定种子 \code{0x5245534F55524345} 不是随机演示。参考模型独立保存每个
资源是否活动、动作登记顺序和预期终态；被测实现只暴露公共接口。每轮随机选择
登记、提交、显式回滚、容量耗尽或动作失败，并比较调用次数、顺序、状态与
资源归零。失败会报告种子和精确迭代，可以稳定重现并沉淀为单元样例。

组合测试不伪造统计。它初始化真实 2-bit frame state、buddy、四级页表、
32 TiB KVA 描述符和栈管理器，创建四个双 guard 栈。活动快照必须观测
4 个栈、16 个物理数据页、24 个 KVA 页、16 个 present 数据页和 8 个
not-present guard；逆序销毁后差异掩码必须为零。

具名 256 MiB functional smoke 运行完整 Shell、IPC、文件系统、用户故障隔离
和资源快照，不是只看 \code{READY} 的简化镜像。64 GiB capacity 使用同一
代码，另要求在 4 GiB 以上真实申请、写回和回收页帧。64 MiB bootstrap 保留
启动与历史故障路径。三种 RAM 规格只改变 QEMU 参数，不通过条件编译切回固定
资源表。

\subsection{可观测性与阶段完成边界}

引用获取、补偿动作、逐页 map/unmap 都是热路径，不逐次输出串口。日志只在
整个事务提交或完整回收后输出稳定标记；QEMU runner 对 tracked field、
changed mask 和通过标记做精确次数检查，并继续给每行附加宿主单调相对时间。
来宾时间仍来自 PIT，两种时间域不会混在一个字段中。

v1.1 完成时共有 97 项 CTest。新增五项分别是原语单元测试、十万事务随机
测试、快照单元测试、真实跨层生命周期集成测试和 256 MiB QEMU functional
smoke。全部历史复位、长模式、异常、持久化和故障镜像继续保留，因此“新功能
通过”不能掩盖旧边界回归。

v1.1 完成时的受控边界也必须明确：引用计数尚未接入动态 KernelObject，
未来六个快照字段仍为零，\code{ScopeRollback} 的动作是同步函数且没有锁。
v1.2 随后把 Process/Thread 当前量接入既有快照；v1.4 又以管理器锁、强引用
和类型 finalizer 完成动态对象的最后释放，但仍不提供弱引用。先建立可证明的
单核语义，再扩展并发，不用一个未经定义的 atomic 字段制造“已经线程安全”
的错觉。

\section{v0.6 高半区早期堆用最小策略验证对象分配}

heap 虚拟基址为 \code{0xFFFF800000000000}，容量 64 KiB。页表管理器为
16 个虚拟页分别取得物理帧；这些帧不要求连续，因为虚拟地址负责把它们组织成
连续区间。每页权限为 RW/NX。

v0.6 的 \code{KernelHeap} 保存 base、size、consumed 和 allocation count。
分配大小必须非零，对齐必须是二的幂。当前地址先检查
\code{base+consumed} 溢出，再计算：
\[
  a=(current+alignment-1)\mathbin{\&}\sim(alignment-1).
\]
padding 和请求大小分别与剩余容量比较，避免组合加法先溢出。只有全部检查
通过才更新 consumed、次数和输出指针；失败调用不会把调用者指针改成半有效值。

这个 heap 故意没有 \code{free}。启动期的页表、设备对象和全局服务通常与
内核同寿命，单调策略没有碎片合并、双重释放和锁状态。它的缺点同样明确：
容量一旦消耗不能回收，所以后续通用堆仍需 size class、slab 或其他所有权
策略，不能把早期 heap 当成最终实现。

v0.6 目标自检先请求 64 字节、16 字节对齐对象，再请求 4096 字节、
4096 字节对齐
对象，分别写入
\code{0x13579BDF2468ACE0} 和 \code{0xC001D00DC0FFEE11} 后读回。
\code{HEAP\_SELF\_TEST\_PASSED} 因而同时证明高半区地址 canonical、四级
映射存在、页可写、NX/WP 没有误伤数据、对齐算法和真实访存都成立。

未来可释放分配器不能在持有自身锁时调用可能再次分配的日志格式化。panic 和
低内存路径继续使用预分配缓冲与串口最小输出，避免诊断代码递归触发原故障。

\section{v1.1 用边界标记把早期堆推进到完整生命周期}

\subsection{为什么“能够分配”还不是资源管理}

单调堆的 consumed 指针只回答“下一块从哪里开始”。它没有记录一个返回指针
对应多大的物理块，也没有办法判断相邻空间是否空闲，因此不存在可靠的
\code{free}。这不是少写一个接口的问题，而是缺少恢复所有权所需的信息。

动态 Thread、内核栈、FileTable 和缓存对象都有明确销毁点。若仍把它们放进
单调堆，功能测试可能短时间成功，长时间创建/退出却只会让 consumed 单向增长。
这类泄漏尤其危险：它不一定破坏当前对象，而是在距离根因很远的后来请求上
表现为耗尽。v1.1 因而分别建立通用字节块和 buddy 页块的申请、释放与守恒，
再在其上推进 KVA 与类型缓存；每层只回答一种所有权问题。

\subsection{块头同时服务物理邻接和空闲集合}

64 KiB 后备虚拟区间保持不变，块头固定为 48 字节并按 16 字节对齐：

\begin{center}
\begin{tabular}{lll}
\toprule
字段 & 宽度 & 作用\\
\midrule
\code{size_bytes} & 64 bit & 当前物理块的完整跨度\\
\code{previous_block_size_bytes} & 64 bit & 从当前头反向定位前块\\
\code{requested_size_bytes} & 64 bit & 区分调用者负载与内部开销\\
\code{state_signature} & 64 bit & 区分 free、allocated 与旧 released 头\\
\code{previous_free_block} & 64 bit & 空闲链反向链接\\
\code{next_free_block} & 64 bit & 空闲链正向链接\\
\bottomrule
\end{tabular}
\end{center}

物理块链与空闲链不是同一件事。前者必须无缝覆盖整个 heap，包括活动块；
后者只索引 free 块，并按地址递增。对地址为 \(b_i\)、长度为 \(s_i\) 的块，
相邻关系为
\[
  b_{i+1}=b_i+s_i,\qquad
  b_{i-1}=b_i-\operatorname{previous_size}(b_i).
\]
保存前块长度就是 boundary tag 的核心。释放时可以在常数时间找到两个物理
邻居，无需从 heap 起点重新搜索；空闲链则让分配只遍历可用候选。

48 字节头看起来昂贵，这是本阶段有意识的取舍。一个 16 字节小对象会占用至少
64 字节块，空间效率不高；但每个字段都可在 GDB 和一致性检查器中直接解释。
后续固定尺寸热对象交给 type cache，通用 heap 不需要同时承担 slab 的职责。

\subsection{对齐会同时产生前缀、活动块和后缀}

请求负载长度 \(n\)、对齐 \(a\) 时，先取
\[
  a_{\mathrm{effective}}=\max(a,16).
\]
负载必须紧跟 48 字节块头，因此先对
\(\operatorname{block_begin}+48\) 向上对齐，再反推出活动块起点。若前缀
非零但小于最小块 64 字节，它既不能进入空闲链，又不能放在负载与块头之间；
实现继续前进一个对齐单位，直到前缀为零或足够形成完整块。

活动块初始跨度为
\[
  \operatorname{align_up}(48+n,16).
\]
尾部不足 64 字节时并入活动块，避免制造永远不能分配的碎片；否则尾部形成
新的 free 块。只有候选地址、三段长度、溢出和边界全部验证后，best-fit
选择才会提交。失败调用不修改空闲链、块头、统计或调用者输出指针。

best-fit 从所有可容纳请求的 free 块中选最小者。它仍是线性搜索，也不保证
消灭外部碎片；v1.1 此时的 heap 只有 64 KiB，这个策略的可解释性和确定性比
复杂索引更重要。对象尺寸分布明确后再用 type cache 消除高频扫描与块头开销。

\subsection{释放先预检，再完成唯一一次状态提交}

\code{TryRelease} 只接受原分配返回的精确负载首地址。它先验证：

\begin{enumerate}[label=\arabic*.]
  \item 指针非空、位于 heap 内并满足最小对齐；
  \item 负载前 48 字节确实是 allocated 块头；
  \item 该块确实出现在从 heap 基址遍历得到的物理块链中；
  \item requested、size 和活动统计不会下溢；
  \item 前后块边界、状态和反向长度彼此一致。
\end{enumerate}

全部预检成功后，当前块才改为 released。若前块 free，先从空闲链摘除并向前
扩展；若后块 free，再摘除并向后扩展。最后更新后继块保存的 previous size，
把一个合并结果按地址插回空闲链并提交统计。于是

\[
  \mathrm{successful_allocation_count}
  =
  \mathrm{releases}
  +
  \mathrm{allocation_count}
\]

在任意合法时刻都成立。重复释放会命中 free 或 released 签名并明确失败；
内部指针即使碰巧按 16 字节对齐，也无法通过物理块成员检查。

\subsection{一致性检查不能只数链表节点}

两个集合节点数相同，不代表它们包含相同块。\code{Validate} 先沿物理块链
验证尺寸、对齐、状态、前块长度和完整覆盖；再沿空闲链验证地址递增、双向链接
和有界无环；最后逐个确认每个物理 free 块在空闲链中恰好出现一次。它还重新
计算活动请求字节、活动块完整跨度和活动数量，与保存统计交叉比较。

统计特意区分 requested 与 consumed。前者是调用者可见负载，后者包含 48 字节
块头和被吸收的短尾部。\code{remaining_bytes} 只是
\(\mathrm{capacity}-\mathrm{consumed}\)，可能散落在多个块；真正的单次
上限由 \code{largest_free_allocation_bytes} 描述。峰值则记录启动周期曾经
占用的最大完整跨度。

\subsection{测试从样例释放扩展到十万步状态机}

单元测试覆盖无效区间、重复初始化、五种对齐、耗尽失败原子性、空/区间外/
内部/重复释放、乱序三向合并、起始地址复用和统计守恒。集成测试使用目标同样
的 64 KiB 区间，混合分配小对象、缓存行对象与 4 KiB 对齐对象，写入模式后
乱序释放，再申请整个最大连续负载。

固定种子 \code{0x4845415056313031} 执行 100000 次随机申请与释放。参考
模型独立保存每个活动负载的地址、长度、对齐与首尾模式；每一步核对数据未被
邻接分裂破坏、区间不重叠、活动请求字节和累计计数一致，并周期性执行完整
拓扑检查。失败输出包含种子与精确迭代号。

QEMU 自检仍完成 64 字节/16 字节对齐和 4096 字节/4096 字节对齐的真实
写回，但随后逆序释放并执行 \code{Validate}。日志要求
\code{HEAP\_ACTIVE\_ALLOCATIONS=0}，同时报告非零峰值和合并后的最大连续
空闲负载，最后才输出 \code{HEAP\_SELF\_TEST\_PASSED}。分裂和释放不逐次
写串口，避免诊断改变时序。

v1.1 实现仍有清晰边界：后备区当时固定为 64 KiB，没有内部并发锁，也不能在
IRQ、NMI 或 panic 中调用。v1.4 为动态对象和文件表把后备区扩为 512 KiB，
但仍不是按页增长 heap。buddy 与 KVA 已由前文的独立页层和虚拟区间层承担；
动态 guard stack 与空中间页表回收均已由独立模块闭环，按页增长仍是后续
能力，而不是 heap 接口暗中承担的职责。固定尺寸 type cache 已经作为下一层
独立机制建立，下面沿其历史、布局和状态机继续展开。

\section{v1.1 用固定尺寸类型缓存分离对象容量}

\subsection{从任意字节块到同构对象}

通用堆接受任意尺寸和二次幂对齐。这样的接口必须为每次请求回答三个问题：
哪个空闲块容得下、对齐会怎样切割前缀、剩余后缀是否还能成为独立块。即使连续
申请的都是 64 字节 Thread 对象，堆也不能从接口上知道下一次请求仍然同型。

内核中的许多资源却天然同构。一个特定版本的 Thread、VMA、等待节点或打开
文件描述具有固定 \code{sizeof} 与 \code{alignof}；它们会大量创建、释放，
又需要独立的容量、峰值和泄漏证据。若每个对象直接走 48 字节块头的通用堆，
64 字节负载至少要消耗 112 字节左右，还要在空闲链中搜索。对象本身很小，
通用性开销反而接近负载。

20 世纪 80 年代末的 SunOS 内核把这种观察发展成 slab 分配思想：为一种对象
准备一批同尺寸槽，空闲对象保留在缓存中，后续创建可以复用已经对齐的存储。
Linux 后来的 SLAB、SLUB 等实现继续加入多个 slab 的空/部分/满状态、
per-CPU 快速路径、NUMA、poison、red-zone 和内存压力回收。名称与内部数据
结构会变化，核心分层却稳定：

\begin{center}
\begin{tabular}{lll}
\toprule
层 & 管理单位 & 回答的问题\\
\midrule
buddy & \(2^k\) 个连续物理页 & 哪些 PFN 由谁持有\\
KVA/页表 & 虚拟区间与映射 & 页出现在哪个地址、具有什么权限\\
通用 heap & 任意字节块 & 一个已映射区间怎样分裂与合并\\
type cache & 同尺寸对象槽 & 某一类型还有多少可交付对象\\
\bottomrule
\end{tabular}
\end{center}

本项目不能因为最终方向类似 SLUB，就一次复制其全部复杂度。type cache 增量
落地时 KVA 与 Thread/IRQ 锁语义都尚未冻结，因此没有同时引入多 slab 增长和
per-CPU 缓存，避免把页后备、地址空间、并发与对象析构错误混在一起。KVA 现已
由前文的独立闭环完成；缓存仍保持一个固定容量、一个 heap 后备块、可以整体
销毁的既有协议，等待按页 slab 作为单独增量接入。

\subsection{原始核心与模板包装为什么分成两层}

\code{KernelFixedObjectCache} 接受对象字节数、对象对齐和容量，返回无类型
存储；\code{KernelTypeCache<ObjectType>} 只负责把
\code{sizeof(ObjectType)} 与 \code{alignof(ObjectType)} 送给核心，并把
成功结果转回 \code{ObjectType*}。布局、位图、链表、统计和错误路径只有一份
实现。

这层包装不调用构造函数，也不在释放时调用析构函数。原因不是 C++ 对象生命期
不重要，恰恰是它必须由所有者明确决定。缓存只知道“一段适合放置 T 的存储”
何时空闲，不知道 Thread 的等待队列是否已摘除、FileDescription 的引用是否
归零，也不知道一次构造失败需要逆序释放哪些资源。把这些动作偷偷放进通用
模板，会让真正的所有权关系消失。

当前目标自检对象只包含八个 \code{uint64_t}，属于可直接写入的平凡存储。
未来非平凡 KernelObject 必须在所属模块中明确形成：

\begin{lstlisting}[language=C++]
storage = cache.TryAcquire()
construct object and acquire child resources
if any step fails:
    roll back committed child resources
    return storage to cache

on final reference:
    detach object from all owners
    destroy object state
    return storage to cache
\end{lstlisting}

因此“模板返回类型化指针”不等于“缓存拥有对象语义”。接口注释把这条边界
写在声明旁，ELF 审计继续拒绝隐藏构造器和运行时符号。

\subsection{一个后备块怎样同时容纳位图与对象}

设对象大小为 \(Z\)，语言要求对齐为 \(A_T\)，容量为 \(N\)。空闲槽首要保存
一个 64 位下一索引，所以内部存储对齐与槽步长为：

\[
\begin{aligned}
 B &= \left\lceil\frac{N}{8}\right\rceil,\\
 A &= \max(A_T,8),\\
 S &= \operatorname{align\_up}(\max(Z,8),A),\\
 O &= \operatorname{align\_up}(B,A),\\
 H &= O+N S.
\end{aligned}
\]

\(B\) 是活动位图字节数，\(A\) 是后备申请所需对齐，\(S\) 是槽步长，
\(O\) 是对象区相对后备块起点的偏移，\(H\) 是唯一 heap 请求长度。布局为：

\begin{verbatim}
backing allocation, H bytes
low address
┌──────────────┬───────────────┬────────┬────────┬─────┬──────────┐
│ bitmap B byte│ align padding │ slot 0 │ slot 1 │ ... │ slot N-1 │
└──────────────┴───────────────┴────────┴────────┴─────┴──────────┘
                                  each slot S bytes, aligned to A
high address
\end{verbatim}

位图与槽共用一个后备块有两项可验证结果。第一，缓存初始化成功只会让 heap
活动申请数增加一；第二，活动对象归零后一次 \code{Destroy} 就能把全部缓存
元数据和槽归还，不存在忘记释放旁路位图的第二条路径。

每个算式都可能溢出。\code{align_up(value,alignment)} 不能先无条件执行
\code{value+alignment-1}；实现先验证 value 不大于
\code{UINT64_MAX-(alignment-1)}。乘法先检查 \(N\le
\mathrm{UINT64\_MAX}/S\)，最后再检查 \(O\le\mathrm{UINT64\_MAX}-NS\)。
这些检查发生在调用 heap 之前，所以非法配置不能改变全局堆累计值。

\subsection{32 个缓存行对象为什么得到 \code{0x840}}

目标自检选择 \(Z=64\)、\(A_T=64\)、\(N=32\)。代入：

\[
\begin{aligned}
B&=\lceil32/8\rceil=4,\\
A&=64,\\
S&=\operatorname{align\_up}(64,64)=64,\\
O&=\operatorname{align\_up}(4,64)=64,\\
H&=64+32\times64=2112=\mathtt{0x840}.
\end{aligned}
\]

4 字节位图之后有 60 字节填充。这看似浪费，却确保第一个槽和随后每个
64 字节步长槽都从缓存行边界开始。这里的“缓存行”只是一种代表性对齐，
项目没有把 64 字节写成所有 x86-64 CPU 的架构 ABI；真实处理器的缓存拓扑
可以由 CPUID 单独发现。该自检只要求请求的语言对齐被分配器兑现。

单元测试另用 \(Z=1\)、\(A_T=1\)、\(N=9\) 验证相反边界：

\[
B=2,\quad A=8,\quad S=8,\quad O=8,\quad H=8+9\times8=80.
\]

容量 9 迫使位图使用第二字节的一个有效 bit，其余七个尾 bit 必须保持零；
1 字节对象又迫使槽扩展到 8 字节以保存空闲索引。这不是为了支持某个实际
1 字节内核对象，而是把两个容易被普通 32 槽样例掩盖的布局分支变成证据。

\subsection{活动位图与槽内空闲链各自证明什么}

每个槽有且只有两种状态：

\begin{longtable}{L{0.14\textwidth}L{0.22\textwidth}L{0.50\textwidth}}
\toprule
\textbf{bitmap} & \textbf{所有者} & \textbf{槽内容解释}\\
\midrule
\endfirsthead
\toprule
\textbf{bitmap} & \textbf{所有者} & \textbf{槽内容解释}\\
\midrule
\endhead
0 & cache & 槽首 \code{uint64_t} 是下一空闲槽索引\\
1 & caller & 全部 \(S\) 字节属于调用者，缓存不得写入\\
\bottomrule
\end{longtable}

只保存空闲链不够。活动对象正在使用槽首，缓存不能在那里留下链节点；重复
释放若没有位图，会把同一槽插入两次，之后两次申请可能向两个调用者交付同一
地址。只保存位图也可以辨别活动，但申请最坏要线性扫描 \(N\) 个 bit。释放
时已经知道槽索引，把它压入链头即可让下一次申请为 \(O(1)\)。

位图定位采用：
\[
  \mathrm{byte}=i/8,\qquad
  \mathrm{bit}=i\mathbin{\&}7.
\]
链尾使用 \code{UINT64_MAX}，有效索引严格小于容量。初始化把槽
\(i\) 的 next 写成 \(i+1\)，最后一槽写链尾，链头设为 0。由于对象区和
步长至少按 8 字节对齐，这些 64 位读写在 x86-64 上不会产生未对齐访问。

空闲链采用 LIFO，而不是保持地址排序。释放只需三步：写旧 head、更新 head、
增加 free；最近释放的存储也更可能仍在 CPU 数据缓存中。LIFO 不是公平性
契约，调用者不能依赖长期地址顺序；目标自检只用“刚完成一组确定释放后应复用
最后槽”证明链提交方向正确。

\subsection{初始化是一项先申请、后发布的事务}

\code{Initialize} 依次执行：

\begin{enumerate}[label=\arabic*.]
  \item 拒绝重复初始化、零对象尺寸、非二次幂对齐和零容量；
  \item 在局部 \code{KernelTypeCacheLayout} 中计算全部尺寸与溢出；
  \item 向 heap 申请一个 \(H\) 字节、\(A\) 对齐的后备块；
  \item 只有 heap 成功后才发布 heap 指针、位图、对象基址和配置字段；
  \item 清零 \(B\) 个位图字节，写入 \(N\) 个空闲 next 索引；
  \item 设置 active=0、free=\(N\)、累计值和峰值为零。
\end{enumerate}

若前三步失败，缓存仍处于 \code{NotInitialized}；不会出现“保存了 heap
指针却没有后备块”的半状态。heap 自身也保证分配失败不改输出和拓扑，于是
两层事务边界可以组合。

初始化不会清零全部对象槽。空闲链本来就要占用每槽首 8 字节，而其余字节没有
安全语义；对象所有者必须初始化自己将读取的字段。无条件清零不仅掩盖遗漏的
对象初始化，还会让缓存容量越大、启动时间越长，却没有改变所有权。

\subsection{申请路径在写位图前验证局部守恒}

\code{TryAcquire} 先核对：

\[
\begin{aligned}
\mathrm{active}+\mathrm{free}&=N,\\
\mathrm{successful}-\mathrm{releases}&=\mathrm{active},\\
\mathrm{active}&\le\mathrm{peak}\le N.
\end{aligned}
\]

若 free=0，head 必须恰为链尾，然后返回 \code{OutOfObjects}。输出引用保留
调用者原值，因此调用者不能把耗尽误当成一个新的空指针对象。若仍有空闲槽，
head 必须在范围内且位图为 0；下一索引必须与剩余 free 数相符、不能指向自身
或一个活动槽。累计申请已为最大 64 位值时返回 \code{CounterOverflow}，
不能回绕为零。

提交顺序为：

\begin{lstlisting}[language=C++]
index = free_head
next = slot[index].next
bitmap[index] = active
free_head = next
active += 1
free -= 1
successful += 1
peak = max(peak, active)
output = address_of(index)
\end{lstlisting}

输出最后写入很重要。之前任一验证失败都不应把一个尚未拥有的槽地址暴露给
调用者。当前是串行路径，所以这些普通读写不使用原子指令；这不是宣称编译器
或 CPU 会替未来并发保持事务。

\subsection{释放先从地址恢复索引，再接触槽内容}

释放路径不能先把 next 写进对象，因为输入指针可能位于另一个对象内部。实现
先把指针解释为 64 位地址，并验证它落在：
\[
  [\mathrm{object\_base},\,
    \mathrm{object\_base}+N S).
\]
随后计算 offset，要求 \(\mathrm{offset}\bmod S=0\)，再取
\(i=\mathrm{offset}/S\)。只有位图 \(i\) 为 1，才说明这是当前活动槽的精确
首地址。

空指针返回 \code{NullObject}；区间外或槽内指针返回
\code{InvalidObject}；位图已经为 0 返回 \code{ObjectNotActive}。三者分开
使调试器能区分调用方没有对象、地址算错和生命周期重复结束。验证计数与当前
free head 后，提交为：

\begin{lstlisting}[language=C++]
bitmap[index] = free
slot[index].next = free_head
free_head = index
active -= 1
free += 1
releases += 1
\end{lstlisting}

清位和写 next 都发生在全部拒绝条件之后，因此重复释放不能把现有链截断。
释放成功之后，旧对象的槽首已经是缓存元数据；调用者继续读取它属于 use-after-
release，而不是缓存应保持的数据契约。

\subsection{销毁为什么必须先证明没有活动对象}

缓存后备块同时包含所有槽。只要一个对象仍活动，归还该块就会让其指针指向
heap 可能重新分配的区域，形成跨类型 use-after-free。因此
\code{Destroy} 在 active 非零时返回 \code{ActiveObjectsRemain}，缓存保持
原状。

active 为零仍不足以立即释放。空闲链可能因更早的错误断裂，位图尾部也可能
损坏；\code{Destroy} 先调用完整 \code{Validate}。只有校验成功且 heap
\code{TryRelease(backing)} 也成功，才清空内部指针和统计。若 heap 拒绝，
返回 \code{BackingReleaseFailed} 并保留缓存状态，调试器仍能检查原现场。

销毁前由调用者取得 \code{Statistics} 快照，因为成功销毁后对象已经回到
未初始化状态，旧统计不再描述一个存在的缓存。目标内存管理器正是先保存
33/33/32 的自检快照，再销毁并记录到全局启动统计。

\subsection{完整校验怎样在不分配临时内存时发现环}

\code{Validate} 首先根据保存的 \(Z,A_T,N\) 重算 \(B,A,S,O,H\)，核对
位图确实等于后备块起点、对象基址等于 backing+\(O\)，再验证全部计数关系。
若 \(N\) 不是 8 的倍数，最后一个位图字节中超出容量的高 bit 必须为零。

然后扫描 \(0..N-1\) 的位图，得到 observed active。空闲链从 head 开始，
每访问一个索引就要求：

\begin{itemize}
  \item 已访问步数小于 \(N\)；
  \item 索引小于 \(N\)；
  \item 对应位图为 0。
\end{itemize}

单向链中若再次访问某节点，之后路径必然重复，步数最终达到容量；因而不需要
另分配 visited 位图也能在下一次解引用前识别环。断链或未连入的空闲槽则让
observed free 小于保存的 free。最后要求 observed active/free 分别等于
统计，两者之和为容量。

校验器不能为了验证缓存再从同一缓存申请临时节点，也不应调用可能依赖 heap 的
日志格式化。它只返回状态；冷路径启动代码在整个内存初始化成功后统一打印
里程碑。

\subsection{三层宿主测试分别攻击哪一条边界}

单元测试从 API 失败语义开始：未初始化申请/释放/销毁、无效配置、尺寸溢出、
heap 后备不足、重复初始化、全容量耗尽、输出保持、空/外部/内部指针、活动
对象销毁、重复释放和 LIFO 复用。最后把全部槽归还，要求唯一 heap 后备块
消失。前述 1 字节、9 槽配置专门覆盖最小步长与尾 bit。

集成测试在一个 64 KiB heap 上同时建立三种缓存：

\begin{center}
\begin{tabular}{rrrr}
\toprule
对象 & 对齐 & 容量 & 关注点\\
\midrule
SmallObject & 8 & 64 & 普通小对象与交错释放\\
CacheLineObject & 64 & 32 & 缓存行对齐与数据保持\\
LargeObject & 256 & 8 & 大对齐和较大槽\\
\bottomrule
\end{tabular}
\end{center}

测试填满三者，释放 small/large 的偶数槽，同时完整读回仍活动的 cache-line
模式。带活动对象的缓存销毁必须失败；随后释放奇数槽和全部 cache-line 槽，
按不同于创建顺序销毁三者。heap 最终必须恢复单一最大空闲负载，累计三次后备
申请等于三次释放。

固定种子 \code{0x5459504543414348} 执行 100000 步。独立记录保存每个活动
指针及首尾模式，随机选择申请或某个活动序号释放；每步核对活动地址唯一、
数据未被空闲链覆盖、active/free 与累计值一致。每 257 步执行完整校验，
每 509 步的相应释放机会再尝试一次重复释放。失败报告包含种子与迭代号，
不会留下只能“再跑一次看看”的偶现结果。

\subsection{QEMU 自检把软件状态机连接到真实映射}

宿主模型证明算法，却不能证明高半区地址真的由当前 CR3 映射，也不能证明
freestanding 目标生成的 64 位加载/存储能访问后备页。目标自检在 heap 自检
完成后执行：

\begin{enumerate}[label=\arabic*.]
  \item 建立 32 槽、64 字节对象、64 字节对齐的缓存；
  \item 申请全部槽，逐地址验证低 6 bit 为零，在索引 0 与 7 写入带槽序号的
        两组 64 位模式；
  \item 第 33 次申请必须耗尽且保留输出，随后读回全部 64 组边界模式；
  \item 先释放 0,2,\ldots,30，重复释放槽 0 必须被拒绝；
  \item 再释放 1,3,\ldots,31，此时 LIFO head 应为槽 31；
  \item 再申请一次必须得到原槽 31，然后立即归还；
  \item 完整校验后保存统计，销毁缓存并再次校验 heap。
\end{enumerate}

自检进入前 heap 已完成两次自身申请/释放。缓存销毁后，当前 consumed、
allocation count、active requested 和 largest free 必须与进入前相同；
累计 heap 申请与释放各增加一次。只有这一整条链成立，内核才在冷路径输出：

\begin{verbatim}
[OS][KERNEL] TYPE_CACHE_READY
[OS][KERNEL] TYPE_CACHE_OBJECT_SIZE_BYTES=0x0000000000000040
[OS][KERNEL] TYPE_CACHE_OBJECT_ALIGNMENT_BYTES=0x0000000000000040
[OS][KERNEL] TYPE_CACHE_SLOT_STRIDE_BYTES=0x0000000000000040
[OS][KERNEL] TYPE_CACHE_CAPACITY=0x0000000000000020
[OS][KERNEL] TYPE_CACHE_BACKING_STORAGE_BYTES=0x0000000000000840
[OS][KERNEL] TYPE_CACHE_ACTIVE_OBJECTS=0x0000000000000000
[OS][KERNEL] TYPE_CACHE_FREE_OBJECTS=0x0000000000000020
[OS][KERNEL] TYPE_CACHE_SUCCESSFUL_ALLOCATIONS=0x0000000000000021
[OS][KERNEL] TYPE_CACHE_RELEASES=0x0000000000000021
[OS][KERNEL] TYPE_CACHE_PEAK_ACTIVE_OBJECTS=0x0000000000000020
[OS][KERNEL] TYPE_CACHE_SELF_TEST_PASSED
\end{verbatim}

每次申请或释放都打印会产生两个反效果：115200 波特串口把分配器变成 I/O
基准，日志路径还可能在未来再次申请内存。当前只更新固定宽度计数，在事务
完成后打印一次，使日志既能证明容量和守恒，又不反向改变热路径。

\subsection{为什么当前没有原子指令和内部锁}

x86-64 对齐的普通 64 位 load/store 足以让当前 BSP 串行启动代码读写索引，
但它不构成多个执行流之间的互斥。若 Thread A 与 IRQ 同时弹出同一个 head，
两者都可能看到旧索引并交付同一槽；仅把 head 改成 atomic 也不能原子提交
位图和三项计数。

v1.2 将冻结三种调用环境：普通 Thread 上下文、与本 CPU IRQ 共享的短临界区、
以及绝不分配的 NMI/panic。具体缓存的所有者届时用 \code{SpinLock} 或
\code{IrqSaveSpinLock} 包围完整状态转换，锁顺序由对象模块声明。当前缓存
内部不擅自 \code{CLI}，也不加入一个尚无调用上下文语义的锁；IRQ、NMI 和
panic 明确不得调用。

未来 per-CPU magazine 可以让常见申请不触及共享锁，但那是性能层。即使快速
路径改变，重复释放拒绝、容量守恒、销毁前 active=0 和后备所有权仍必须与
本次状态机一致。

\subsection{从固定单块到多 slab 的演进边界}

这里的固定缓存容量初始化后不变。它没有空/部分/满 slab 列表，也不能只归还
一部分后备页；在 v1.1 时后备仍受 64 KiB heap 总容量约束，v1.4 的通用 heap
虽已扩为 512 KiB，缓存自身仍是单后备块。这足以验证固定槽所有权，却不是
大规模内核的最终分配策略。

KVA 已经能够从高半区分配虚拟区间，并与 buddy、页表组成经过整机验证的跨层
事务。下一步让类型缓存拥有多个按页后备 slab；空 slab 才能撤销映射、归还
物理页和 KVA。动态内核栈使用同一 KVA/页后备事务，但额外长期保留
not-present guard。已经落地的页表回收保证撤销映射不会永久泄漏独占
中间表。

v1.2 的 Process/Thread 对等路径通过后，固定 PCB 才迁入动态缓存；v1.4
进一步把引用计数 KernelObject 和动态 FileDescription 建立在通用可回收
heap 上。类型缓存继续服务尺寸完全一致的热对象，不被强行用于带私有头和异构
载荷的对象。这一顺序保留一条重要工程原则：基础设施可以先存在，旧生产路径
必须在新对象模型拥有对等测试之前继续工作，不能为了“看起来已经动态化”制造不可启动的
中间版本。

\section{内存测试从纯函数扩展到处理器执行}

四层验证各自回答不同问题：

\begin{enumerate}[label=\arabic*.]
  \item 单元测试验证最高可用 RAM 与高端保留洞的区别、元数据区搜索、
        64 GiB 状态尺寸、高地址范围分配、direct-map canonical 边界、
        heap 原子失败、释放合并和表项权限编码；
  \item 集成测试同时保留 64 MiB 启动图，并构造 3--4 GiB 洞和总计 64 GiB
        可用 RAM，核对 16777216 个 free frame 和 4 GiB 以上分配/回收；
  \item 固定种子 \code{0x6D656D6F72793634} 完成 8192 组叶表项往返，
        执行 4096 步随机 allocate/release，再执行 1024 轮随机高地址窗口
        分配并立即回收；独立 heap、type cache 与 KVA 种子分别再执行
        100000 步字节块、固定槽和虚拟区间申请/释放；
  \item QEMU 主路径固定使用 64 GiB RAM，真实加载 CR3、访问 heap，并要求
        direct-map 精确覆盖 64 GiB、使用 32768 个 2 MiB 页，在
        \code{0x0000000100001000} 或更高帧写入读回两个 64 位模式；失败路径
        继续证明 \code{MEMORY\_MAP\_INVALID} 和错误码 3 的写保护 \#PF。
\end{enumerate}

随机测试不是用随机数替代设计。种子、迭代号和性质名都写入失败报告，任何失败
可以精确重放。参考模型只记录“这一帧是否 allocated”，不调用被测分配器的
内部状态，因此不会与实现犯同一种编码错误。

\section{panic 是终止协议而不是普通日志函数}

panic 首先禁止本地中断并阻止其他核心继续修改共享状态，随后输出原因、CPU、
RIP、寄存器、控制寄存器和有限栈回溯。输出函数不得动态分配或获取未知锁。
最后进入明确停机、调试陷阱或重启策略。

QEMU 测试可让 panic 写入专用退出端口或稳定串口标记，宿主据此区分预期故障
与超时。生产硬件则可能保存崩溃区或触发看门狗，机制边界保持相同。
